\documentclass[a4paper]{scrreprt}
\usepackage[german]{babel}
\usepackage[german]{translator}
\usepackage[utf8]{inputenc}
\usepackage[T1]{fontenc}
\usepackage{blindtext} 
\usepackage{ae}
\usepackage[bookmarks,bookmarksnumbered]{hyperref}
\usepackage{graphicx}
\usepackage{color}
\usepackage[dvipsnames]{xcolor}

\newcounter{Lc}
\newcounter{Hc}
\newcommand{\stepHc}{\stepcounter{Hc}\setcounter{Lc}{0}}
\newcommand{\resetAllCounter}{\setcounter{Lc}{0}\setcounter{Hc}{1}}

%Glossar
\usepackage[nonumberlist, toc, section=chapter, numberedsection=nolabel]{glossaries}
\makenoidxglossaries

%Glossareinträge
\newglossaryentry{Admin}{name=Administrator, description={Ein Administrator oder kurz Admin ist ein Systemverwalter, welcher besondere (weiterführende) Zugriffsrechte hat. Er betreut und verwaltet das System}}
\newglossaryentry{Trainer}{name=Kursleiter, description={Ein Kursleiter ist ein Trainer, welcher einen oder mehrere Kurse leitet}}
\newglossaryentry{Einheiten}{name=Kurseinheiten, description={Jeder Kurs wird in verschiedene Kurseinheiten unterteilt. Dabei entspricht eine Einheit einem stattfindenden Kurstermin}}
\newglossaryentry{Konfiguration}{name=Konfiguration, description={Eine Konfiguration (meist durch den Systemadministrator) beinhaltet die Installation und die Auswahl von Voreinstellungen sowie beispielsweise das Bearbeiten des Impressums}}
\newglossaryentry{Client}{name=Client, description={Ein Client ist ein Programm, welches auf einem Endgerät ausgeführt wird und Dienstleistungen von einem Server in Anspruch nimmt}}
\newglossaryentry{Server}{name=Server, description={Ein Server ist ein Dienstleister, der in einem Netzwerk Daten oder Ressourcen zur Verfügung stellt}}
\newglossaryentry{JSF}{name=JSF, description={engl. Java Server Faces, ist ein Framework zur Entwicklung von grafischen Benutzeroberflächen für Webprogramme.}}
\newglossaryentry{Tomcat}{name=Apache Tomcat, description={Open Source Applikationsserver zum Ausführen von Web-Anwendungen, welche über Java geschrieben wurden}}
\newglossaryentry{PostgreeSQL}{name=PostgreeSQL, description={PostgreeSQL ist ein freies, objektrelationales Datenbank-Verwaltungssystem}}
\newglossaryentry{SMTP-Server}{name=Kursleiter, description={Dies ist ein Simple Mail Transfer Protokoll welches zum Austausch von E-Mails in Computernetzen dient}}
\newglossaryentry{Orgware}{name=Orgware, description={Die Orgware beschreibt Rahmenbedingungen bei IT-Projekten (also alle organisatorischen, methodischen und personellen Maßnahmen), die weder zu der Software noch zu der Hardware gehören}}
\newglossaryentry{Browser}{name=Browser, description={Browser wie Firefox und Internet Explorer sind spezielle Programme zur Websitedarstellung im World Wide Web}}
\newglossaryentry{URL}{name=URL, description={engl. Uniform Resource Locator, bezeichnet die Adresse einer Webseite}}
\newglossaryentry{Id}{name=Identifikationsnummer, description={ist eine Nummer zur eindeutigen Identifikation}}
\newglossaryentry{Default}{name=Default, description={ist eine voreingestellte Eingabevariable}}
\newglossaryentry{SSL}{name=SSL, description={engl. Secure Sockets Layer ist ein Netzwerkprotokoll zur sicheren Datenübertragung im Internet}}
\newglossaryentry{Dummy}{name=Dummy, description={ ist ein Platzhalter}}
\newglossaryentry{JPG}{name=.jpg, description={ Bezeichnung einer Bilddatei}}
\newglossaryentry{CSS}{name=.css, description={ engl. Cascading Style Sheet, wird für die stylistische Anpassung einer Webseite verwendet}}
\newglossaryentry{String}{name=String, description={Ein String ist eine Folge von vordefinierten Zeichen}}
\newglossaryentry{Boolean}{name=Boolean, description={Ein Boolean ist ein Wahrheitswert, also True oder False}}
\newglossaryentry{Integer}{name=Integer, description={Ein Integer ist ein Datentyp zur Speicherung ganzzahliger Werte}}
\newglossaryentry{float}{name=float, description={Ein Float ist ein Datentyp zur Speicherung von Gleitkommazahlen}}
\newglossaryentry{persistent}{name=persistenten, description={persistent bedeutet längerfristig}}
\newglossaryentry{hashen}{name=gehasht, description={bedeutet verschlüssseln}}
\newglossaryentry{Konsistenz}{name=Konsistenz, description={bedeutet Korrektheit}}




\begin{document}
	\thispagestyle{plain}

\begin{titlepage}
    \begin{center}
\begin{figure}[th]
\centering
\includegraphics[width=0.6\linewidth]{logo/name_blau_ofCourse.jpg}
\end{figure}

    	\begin{title}
        	\title{\Huge{\textbf{Kurseinheiten Manager \\ Pflichtenheft\\}}}

		\end{title}
		\hspace{3cm}

        	Software Engineering Praktikum \\
        	Sommersemester 2015\\
        	Universität Passau\\


        	Betreuer: Andreas Stahlbauer\\
        	\hspace{1,5cm}\\
        	Version: 1.0 \\
        	\hspace{1,5cm}\\
        	Datum: 24.04.2015\\[50pt]
        	\textbf{Team 3} \\
            \ \\
    
        
        
        \begin{tabular}{ | l | l | l | l |}
            \hline
            \textbf{Matrikelnummer} & \textbf{Name} & \textbf{Phase} & \textbf{E-Mail}  \\ \hline
            63097 & Katharina Hölzl & Pflichtenheft & hoelzlka@fim.uni-passau.de \\ \hline
            64504 & Ricky Strohmeier& Entwurf & strohric@fim.uni-passau.de  \\ \hline
            64380 & Martin Bachhuber & Feinspezifikation  & bachhube@fim.uni-passau.de \\ \hline
            64080 & Tobias Fuchs & Implementierung  &  fuchstob@fim.uni-passau.de\\ \hline
            61085 & Sebastian Schwarz & Validierung & sebastian@nrschwarz.de \\ \hline  
            58379 & Patrick Cretu  &  Präsentation & cretu@fim.uni-passau.de \\ \hline
        \end{tabular}
        
        \ \\
        \ \\
        \hspace{3 cm}\\
          \textbf{Arbeitspakete Pflichtenheft} \\
          \ \\
        
        \begin{tabular}{ | l | l |}
        	\hline
        	\textbf{Autor} & \textbf{Verantwortungsbereich} \\ \hline
        	 Katharina Hölzl & Globale Testfälle und Testszenarien, Glossar  \\ \hline
        	 Ricky Strohmeier& Zielbestimmung, Qualitätsbestimmung  \\ \hline
        	 Martin Bachhuber & Produkteinsatz, Produktumgebung, Entwicklungsumgebung  \\ \hline
	         Tobias Fuchs & Produktfunktionen \\ \hline
        	 Sebastian Schwarz & Produktdaten, Produktleistungen \\ \hline  
	         Patrick Cretu  &  Benutzeroberfläche, Ergänzung, Benutzerhilfe  \\ \hline
        \end{tabular}
    \end{center}
\end{titlepage}
 
 


% Platzierung des Inhaltsverzeichnisses
\tableofcontents
 
\chapter{Zielbestimmung}
	\begin{tiny}
		RS
	\end{tiny}\\
	   Bei der Webanwendung OfCourse handelt es sich um eine Online-Plattform, welche ein allgemeines Planungssystem für Kurse darstellt. Es ermöglicht Kursleitern und Kursteilnehmern die einfache Absprache von Terminen, was bisher zumeist eine große Herausforderung darstellt.
    \section{Musskriterien}      
    	\subsubsection{Allgemeiner Funktionsumfang:}
      		\begin{itemize}
	      		\item Hilfeseite aufrufen \hyperlink{hilfeSeite}{F10-80}
	      		\item Impressum anzeigen \hyperlink{impressum}{F10-70}
	      		\item Allgemeine Geschäftsbedingungen anzeigen \hyperlink{agb}{F10-90}
      		\end{itemize}
     	\subsubsection{Funktionsumfang für anonymen Benutzer:}
       		\begin{itemize}
	      		\item Durchsuchen aller angebotenen Kurse \hyperlink{kursSuche}{F10-50}
	      		\item Kursanzeige sortieren \hyperlink{kursAnzSort}{F10-40}
	       		\item Registrieren \hyperlink{Registrieren}{F10-20}
       		\end{itemize}
     	\subsubsection{Funktionsumfang für alle registrierten Benutzer:}
			\begin{itemize}
				\item Login \hyperlink{login}{F20-10} und Logout \hyperlink{logout}{F20-90}
				\item Kann sich zu Kursen anmelden \hyperlink{kursAnmelden}{F20-120} und abmelden \hyperlink{kursAbmelden}{F20-130}
				\item Kann sich in Kurseinheiten eintragen \hyperlink{kurseinheitAnmelden}{F20-140} und austragen \hyperlink{kurseinheitAbmelden}{F20-150}
				\item Kann eigene Profildaten editieren \hyperlink{profilEdit}{F20-60}
				\item Kann belegte Kurse anzeigen \hyperlink{kurseAnzeigen}{F20-20}
				\item Einsehen des eigenen Kontostands \hyperlink{kontoAnzeigen}{F20-80} und Geld aufladen \hyperlink{kontoAufladenOn}{F20-180} und \hyperlink{kontoAufladenOff}{F20-190}
				\item Kann Passwort zurücksetzen \hyperlink{passwort}{F20-100}
				\item Funktionsumfang eines Benutzers erweitert sich aufsteigend je nach zugeteilter Rolle
			\end{itemize}
		\subsubsection{Funktionsumfang für Kursleiter:}
			\begin{itemize}
				\item Eigene Kurse bearbeiten \hyperlink{kursEdit}{F30-20}
				\item Kurseinheiten anlegen \hyperlink{kurseinheitAnlegen}{F30-30}, editieren \hyperlink{kurseinheitEdit}{F30-40} und löschen \hyperlink{kurseinheitLoeschen}{F30-70}
				\item Teilnehmer aus Kurs entfernen \hyperlink{kursUserDel}{F30-80}
				\item Nutzer eines Kurses oder einzelner Kurseinheiten benachrichtigen \hyperlink{kursUserMsg}{F30-900}
				\item Accounts bestätigen \hyperlink{addUserKL}{F30-110}
			\end{itemize}
		\subsubsection{Funktionsumfang für Systemadministrator:}
			\begin{itemize}
				\item verwaltet \hyperlink{kursVerwalt}{F40-20}, erstellt \hyperlink{kursErstellen}{F40-10}, bearbeitet \hyperlink{kursBearbeiten}{F40-30} und löscht \hyperlink{kursLoeschen}{F40-40} Kurse und kann zu den Kursen Kursleiter zuordnen \hyperlink{kursKLAdd}{} und wieder entfernen \hyperlink{kursKLDel}{}
				\item Benutzerverwaltung \hyperlink{nutzerVerwalt}{F40-60} einschließlich der Änderung einer Benutzerrolle von Nutzer zu Kursleiter/\gls{Admin} oder \gls{Trainer} zu Nutzer/Administrator \hyperlink{nutzerRolle}{F40-90}
				\item Benutzer suchen \hyperlink{nutzerSuchen}{F40-100}, anlegen \hyperlink{nutzerAnlegen}{F40-50} und löschen \hyperlink{nutzerLoeschen}{F40-80}
				\item Statistiken anzeigen \hyperlink{statistik}{F40-120}
				\item Nutzerguthaben aufladen \hyperlink{guthabenAuf}{F40-130}
				\item System anpassen, beispielsweise eigenes Logo hochladen \hyperlink{logo}{F40-150}, Impressum bearbeiten \hyperlink{impressumBea}{F40-180} oder Registrierungsmodalitäten festlegen \hyperlink{regMod}{F40-160}
			\end{itemize}
			
    \section{Wunschkriterien}
			\begin{itemize}
				\item Englisch als Sprachvariante \hyperlink{spracheAendern}{F10-100}
				\item Terminplaner \hyperlink{terminplaner}{F20-110}
	     		\item Einschränken des angezeigten Kursangebots \hyperlink{kursEinschraenken}{F10-110}
			\end{itemize}
			
	\section{Abgrenzungskriterien}
     		\begin{itemize}
	     		\item Keine Raumplanung/-verwaltung für die Kurse
	     		\item Keine Einsicht in fremde Nutzerprofile, nur Admin kann Nutzerdaten einsehen
	     		\item Kein systeminternes Kommunikationssystem zwischen den Benutzern
	     		\item Keine Funktionen zur erleichterten Bedienung für Nutzer mit Behinderung
	     		\item Keine spezielle Anpassung an Mobilgeräte
     		\end{itemize}
        
  
\chapter{Produkteinsatz}
	\begin{tiny}
		MB
	\end{tiny}\\

    \section{Anwendungsbereiche}
		 Der \gls{Einheiten}-Manager dient der Organisation und Planung von Veranstaltungen. Im System können sich bis zu 80 Benutzer gleichzeitig über verschiedene Veranstaltungen informieren.   
     
	\section{Zielgruppen}
		 \begin{enumerate}
		 	\item Betreiber: Die Webanwendung soll von Betreibern unterschiedlichsten Fachbereichs verwendet werden können, beispielsweise Sportvereine oder Lehrinstitute, welche ihre angebotenen Veranstaltungen und Kurse per Internet organisieren möchten.
		 	\item Benutzer: Es wird zwischen vier verschiedenen Benutzern unterschieden:
		 	\begin{itemize}
		 		\item Anonymer Benutzer: Anonyme Benutzer benötigen normale Kenntnisse im Umgang mit dem Computer. Sie sollten einen \gls{Browser} bedienen können, um die Webseite des \gls{Einheiten}-Manager anzuzeigen. Dort kann man nach bestimmte Veranstaltungen suchen und sich über diese informieren.
		 		\item Registrierter Benutzer: Registrierte Benutzer bedürfen ebenso wie anonyme keine besonderen Vorkenntnisse. Sie sollen einen \gls{Browser} ihrer Wahl nutzen können und benötigen lediglich eine E-Mail-Adresse, um sich im System registrieren und somit ein Benutzerkonto erstellen zu können. Nachfolgend kann nach Veranstaltungen gesucht werden und sich für solche an- oder abmelden.
		 		\item Kursleiter: Ein Kursleiter ist gleichfalls auf wenig besondere IT-Kenntnisse angewiesen. Er sollte allerdings über alle nötigen Informationen seiner Veranstaltung in Kenntnis sein, um jene anlegen und verwalten zu können.
		 		\item \gls{Admin}: Ein \gls{Admin} ist für die Verwaltung und \gls{Konfiguration} des gesamten Systems zuständig. Eine Konfiguration beinhaltet die Installation und die Auswahl von Voreinstellungen, sowie beispielsweise das Bearbeiten des Impressums. Aufgrund dessen werden grundlegende Kenntnisse im Umgang mit HTML, CSS und SQL vorausgesetzt.
		 	\end{itemize}	
		 \end{enumerate}  
        
    
	\section{Betriebsbedingungen}
	       \begin{itemize}
		       	\item Betriebszeit/Betriebsdauer: Ununterbrochener Betrieb mit Ausnahme einer wöchentlichen Wartungszeit von 2 Stunden bzw. 8 Stunden im Monat, welche aber nicht zur Hauptnutzungszeit des Systems abgehalten wird, sofern dies möglich ist.
		       	\item  Im laufenden Betrieb ist keine Aufsicht vonnöten.
	       \end{itemize}
	
			
 
\chapter{Produktumgebung}
	\begin{tiny}
		MB
	\end{tiny}\\
	
	\section{Software}
        \begin{itemize}
      		\item \gls{Client}:
	      		\begin{itemize}
	      			\item Alle gängigen Browser in ihrer jeweiligen aktuellen Version  werden unterstützt: Mozilla Firefox 37, Internet Explorer 11, Google Chrome 42, Apple Safari 8
	      		\end{itemize}
          	\item \gls{Server}:
	            \begin{itemize}
	            	\item Java 8
	            	\item \gls{JSF} 2.2 (Entwicklungsframework)
	            	\item \gls{Tomcat} 8 (Webserver)
	            	\item \gls{PostgreeSQL} 9 (Datenbank)
	            \end{itemize}
        \end{itemize}
        
    \section{Hardware}   
        \begin{itemize}
          	\item Client:
	            \begin{itemize}
	            	\item Endgerät mit Internetzugang
	            \end{itemize}
          	\item Server:
	           \begin{itemize}
		           	\item Rechner mit Internetzugang
		           	\item Applikation/Datenbank-Server: Intel(R) Core(TM) i5-2400 CPU, 8GB
RAM, mindestens 5 GB freier Speicher oder ähnliches
	           \end{itemize}
        \end{itemize}
        
     \section{Orgware}
             \begin{itemize}
               \item Internetverbindung
               \item \gls{Admin}, der die Anwendung einrichtet, die Datenbank wartet
und sichert.
             \end{itemize}

\resetAllCounter
\newcommand{\Func}[1]{\stepcounter{Lc}\textcolor{Blue}{\textbf{/F\arabic{Hc}0-\arabic{Lc}0/} #1} \\}
\newcommand{\FuncW}[1]{\stepcounter{Lc}\textcolor{Green}{\textbf{/F\arabic{Hc}0-\arabic{Lc}0W/} #1} \\}
\newcommand{\FuncBlue}[1]{\textcolor{Blue}{\textbf{#1}}}
\newcommand{\FuncGreen}[1]{\textcolor{Green}{\textbf{#1}}}

\chapter{Produktfunktionen}
\begin{tiny}
	TF\\
\end{tiny}
\ \\
Die Produktfunktionen sind nach Benutzergruppen geordnet. Wunschkriterien sind mit nachgestelltem 'W' gekennzeichnet.
	\section{Anonymer Benutzer}
		\subsection{Grundfunktionen}
			\begin{itemize}
				\item \Func{Aufruf der Startseite.}
					Die Startseite kann mit jedem unterstützten \gls{Browser} (vgl. Produktumgebung -> 3.1 Software -> Client) via \gls{URL} aufgerufen werden.
				\item \Func{Registrierung im System durch Ausfüllen eines Formulars.} \hypertarget{Registrieren}
					Der Nutzer kann sich durch Ausfüllen eines Registrierungsformulars im System registrieren. Die Anwendung generiert für den neu registrierten Nutzer automatisch eine im System eindeutige Identifikationsnummer. Nach Abschicken des Formulars wird an die angegebene E-Mailadresse eine Bestätigungsmail  mit einem Verifizierungslink geschickt, welcher dazu dient, die E-Mailadresse des neuen Benutzers zu bestätigen. Der Nutzer wird dann auf die Startseite des Systems weitergeleitet. Ob sich der Benutzer nun schon in das System einloggen kann, hängt von der Art der Accountaktivierung ab. Diese wird durch den Administrator festgelegt(\FuncBlue{/F40-160/}). Der Account ist entweder durch die E-Mailverifikation automatisch aktiviert worden oder muss manuell durch einen Kursleiter oder einen Administrator aktiviert werden(\FuncBlue{/F30-110/}).
				\item \Func{Anzeigen des Kursangebots.} \hypertarget{kursAngAnz}
					Durch Klicken der Schaltfläche 'Kursangebot' wird dem Nutzer das gesamte Kursangebot angezeigt. Die Anzeige ist automatisch aufsteigend nach dem 'Beginn' der nächsten Kurseinheit des Kurses sortiert. 
				\item \Func{Sortieren des angezeigten Kursangebots.} \hypertarget{kursAnzSort}
					Die angezeigten Kurse können entweder nach 'Titel' des Kurses alphabetisch oder nach dem 'Beginn' der nächsten Kurseinheit des Kurses chronologisch sortiert werden.
				\item \Func{Suche nach Kursen.} \hypertarget{kursSuche}
					Das Kursangebot kann durch Eingabe des Kurstitels und Betätigen der 'Suchen' - Schaltfläche nach einem bestimmten Kurs durchsucht werden.
				\item \Func{Anzeigen der Kursdetails} 
				    Durch Klicken auf einen Kurs des Kursangebots wird der anonyme Nutzer zur 'Kursdetails' - Seite des Kurses weitergeleitet. Auf der Seite werden die Identifikationsnummer des Kurses, der 'Titel', die 'Beschreibung', die 'Kursleiter' und die 'Kurseinheiten' des Kurses angezeigt(\FuncBlue{/D20-10/}). Eine Anmeldung zum Kurs oder zu einer Kurseinheit und die Anzeige der Kursteilnehmer ist für den anonymen Nutzer nicht möglich.
				\item \Func{Anzeige des Impressums.}\hypertarget{impressum}
					Das Impressum kann von jeder Seite des Systems aus angezeigt werden.
				\item \Func{Anzeige der Hilfeseite.} \hypertarget{hilfeSeite}
					Von jeder Seite des Systems kann die Hilfeseite zur entsprechenden Seite des Systems oder die komplette Hilfeseite aufgerufen werden.
				\item \Func{Anzeige der Allgemeinen Geschäftsbedingungen.} \hypertarget{agb}
					Die Allgemeinen Geschäftsbedingungen können von jeder Seite des Systems aus angezeigt werden.
				\item \FuncW{Wechsel der Anzeigesprache.} \hypertarget{spracheAendern}
					Von jeder Seite des Systems aus kann durch Klicken einer entsprechenden Schaltfläche die Anzeigesprache gewechselt werden. Als Default ist 'Deutsch' ausgewählt. Als weitere Sprache ist 'Englisch' vorgesehen. Vom einem Benutzer eingegebene Daten, wie beispielsweise die Beschreibung eines Kurses, werden nicht übersetzt.
				\item \FuncW{Einschränken des angezeigten Kursangebots.} \hypertarget{kursEinschraenken}
					Der Nutzer hat die Möglichkeit das angezeigte Kursangebot auf verschiedene Zeiträume einzuschränken. Vorgesehene Anzeigezeiträume sind 'Tagesangebot' und 'Wochenangebot'. Die Auswahl des Anzeigezeitraums ist über eine Drop-Down-Liste möglich. Bestätigt wird die Auswahl durch Klicken der 'Anzeigen' - Schaltfläche. 
			\end{itemize}

	\section{Registrierter Benutzer}
	Der registrierte Benutzer verfügt über sämtliche Funktionen wie der anonyme Benutzer. Zusätzlich kann er auf folgende Funktionen zurückgreifen.
		\stepHc
		\subsection{Grundfunktionen}
			\begin{itemize}
				\item \Func{Sicheres Einloggen ins System.} \hypertarget{login}
					Durch Eingabe des Benutzernamen, des Passworts und Klicken der 'Anmelden' - Schaltfläche kann sich der Benutzer ins System einloggen. Die Anmeldedaten werden dabei mittels \gls{SSL} übertragen. Nach erfolgreichem Anmelden wird der Nutzer auf die 'Meine Kurse' - Seite weitergeleitet.
				\item \Func{Anzeigen der angemeldeten Kurse.} \hypertarget{kurseAnzeigen}
					Die Kurse zu denen der registrierte Benutzer angemeldet ist werden gesammelt auf der 'Meine Kurse' - Seite angezeigt.
				\item \Func{Detailanzeige der angemeldeten Kurse.}
					Durch Klicken auf einen angemeldeten Kurs werden die Details dieses Kurses angezeigt. Der registrierte Nutzer bekommt die Daten aus \FuncBlue{/F10-60/} angezeigt. Zusätzlich kann er in seinen angemeldeten Kursen die Teilnehmerlisten einsehen und bekommt Details der Kurseinheiten, wie etwa 'Ort' und 'Preis'(vgl. \FuncBlue{/D30-10/}), angezeigt. Außerdem hat der registrierte Benutzer die Möglichkeit sich zu Kurseinheiten an-/abzumelden und sich aus dem Kurs auszutragen.
				\item \Func{Anzeige der Teilnehmer eines angemeldeten Kurses}
				    Durch Betätigen der 'Zu den Teilnehmern' - Schaltfläche auf der 'Kursdetails' - Seite eines angemeldeten Kurses kann der registrierte Nutzer die Teilnehmerliste des Kurses einsehen. 
				\item \Func{Anzeigen der Profilseite.}
					Der registrierte Benutzer besitzt eine Profilseite auf welcher seine Daten gesammelt angezeigt werden(vgl. \FuncBlue{/D40-10/}).
				\item \Func{Ändern der eigenen Benutzerdaten.} \hypertarget{profilEdit}
					Der registrierte Benutzer kann durch Klicken der Schaltfläche 'Bearbeiten' seine eingegebenen Daten bearbeiten. Sollte die E-Mailadresse geändert werden, wird an die neue E-Mailadresse wie beim Registrierungsprozess(\FuncBlue{/F10-20/}) eine Bestätigungsmail mit einem neuen Verifizierungslink geschickt, um die neue E-Mail-Adresse auf Korrektheit zu überprüfen. Der registrierte Nutzer kann alle seine Benutzerdaten ändern mit Ausnahme seines Kontostands, seiner Benutzerrolle und der vom System generierten Identifikationsnummer. Durch Betätigen der 'Speichern' - Schaltfläche werden die Daten gespeichert.
				\item \Func{Hochladen eines Profilbilds.}
					Der registrierte Benutzer kann ein Profilbild hochladen. Das Bild muss das .jpg - Format haben. Anforderungen bezüglich Auflösung und Größe des Bildes werden während der Produktentwicklung festgelegt. Als Default ist ein \gls{Dummy} - Bild gesetzt.
				\item \Func{Anzeigen des aktuellen Kontostands.} \hypertarget{kontoAnzeigen}
					Der registrierte Benutzer kann den aktuellen Kontostand seines Guthabenkontos auf seiner Profilseite einsehen. 
				\item \Func{Aus dem System abmelden mittels 'Abmelden' - Schaltfläche.} \hypertarget{logout}
					Der registrierte Benutzer kann sich  von jeder Seite des Systems aus durch Klicken der Schaltfläche 'Abmelden' aus dem System abmelden und wird auf die Startseite weitergeleitet.
				\item \Func{'Passwort vergessen' - Schaltfläche.} \hypertarget{passwort}
					Klickt der Benutzer die 'Passwort vergessen' - Schaltfläche, so wird er aufgefordert seine im System hinterlegte E-Mailadresse einzugeben. An diese wird dann ein automatisch generiertes Passwort geschickt.
			\end{itemize}	

		\subsection{Kursfunktionen}
			\begin{itemize}
				\item \Func{Detailansicht der Kurse.}
					Der registrierte Nutzer kann sich eine Detailansicht der angebotenen Kurse anzeigen lassen. Angezeigt werden die Daten aus \FuncBlue{/F10-60/}. Außerdem kann sich der Nutzer in der Detailansicht zum Kurs anmelden. Eine Einsicht in die Teilnehmerliste des Kurses ist erst nach der Anmeldung zum Kurs möglich.
				\item \Func{Zu Kurs anmelden.} \hypertarget{kursAnmelden}
					Der registrierte Nutzer kann sich zu beliebig vielen Kursen anmelden. 
					Dafür klickt er in der Detailansicht des gewünschten Kurses auf die Schaltfläche 'Anmelden'.
					Bei der Anmeldung hat der registrierte Nutzer die Möglichkeit sich für Benachrichtigungen zu diesem Kurs zu registrieren(vgl. \FuncBlue{/F20-170/}).
				\item \Func{Von Kurs abmelden.}  \hypertarget{kursAbmelden}
					Der registrierte Nutzer kann sich von Kursen abmelden zu denen er angemeldet ist. Dafür klickt er in der Detailansicht des gewünschten Kurses auf die Schaltfläche 'Abmelden'. Meldet sich der registrierte Benutzer von einem Kurs ab, meldet er sich damit auch automatisch von allen Kurseinheiten ab.
				\item \Func{Zu Kurseinheit anmelden.} \hypertarget{kurseinheitAnmelden}
					Der registrierte Benutzer kann sich zu beliebig vielen Kurseinheiten anmelden. Um
					sich zu Kurseinheiten anmelden zu können, muss der Nutzer bei dem zugehörigen Kurs angemeldet sein. Für die Anmeldung zur Kurseinheit muss der Nutzer in die Detailansicht des Kurses wechseln. Dort findet er die zu diesem Kurs verfügbaren Einheiten und deren Status. Ist eine Kurseinheit noch nicht voll und ist der Nutzer noch nicht zu dieser Einheit angemeldet, kann er sich anmelden. Der Nutzer hat außerdem die Möglichkeit durch Klicken der Schaltfläche 'Alle auswählen' und anschließender Betätigung der Schaltfläche 'Speichern' sich zu allen Einheiten des Kurses auf einmal anzumelden. Vor Betätigen der Schaltfläche 'Speichern' hat der Nutzer aber auch noch die Möglichkeit einzelne Einheiten von der Anmeldung auszuschließen. Voraussetzung um sich zu kostenpflichtigen Kurseinheiten anmelden zu können, ist ein ausreichendes Guthaben auf dem Konto bzw. keine unzulässige Überziehung des Guthabenkontos.
				\item \Func{Von Kurseinheit abmelden.} \hypertarget{kurseinheitAbmelden}
					Der registrierte Benutzer kann sich von Kurseinheiten abmelden, zu denen er angemeldet ist. Um sich von Kurseinheiten eines Kurses abmelden zu können, muss der Nutzer in die Detailansicht des Kurses wechseln. Dort findet er die zu diesem Kurs verfügbaren Einheiten.
					Er kann nun die Einheiten auswählen von denen er sich abmelden möchte und durch Betätigen der Schaltfläche 'Speichern' meldet er sich von den gewählten Einheiten ab. Der für kostenpflichtige Einheiten bezahlte Betrag wird automatisch auf das Konto des registrierten Benutzers zurückgebucht.
					Eine Abmeldung von einer Kurseinheit ist nur bis drei Stunden vor deren Beginn möglich. Sollte nämlich aufgrund der Abmeldung des registrierten Nutzers die minimale Teilnehmerzahl der Kurseinheit unterschritten werden, kann der Kursleiter entscheiden, ob die Kurseinheit stattfindet oder nicht. Durch die Begrenzung der Abmeldefrist hat der Kursleiter somit noch eine gewisse Reaktionszeit und kann gegebenenfalls die anderen Teilnehmer über den Ausfall der Kurseinheit informieren. 
				\item \Func{Nach Kursen suchen - erweitert.}
					Der registrierte Nutzer hat die Möglichkeit das gesamte Kursangebot zu durchsuchen. Über eine Drop-Down-Liste kann der Nutzer einen Suchparameter auswählen. Nach Eingabe eines Suchbegriffs und Klicken der Schaltfläche 'Suchen' wird das Kursangebot nach dem Suchbegriff bzgl. des Suchparameters durchsucht und das Ergebnis angezeigt. 
					Als Suchparameter sind vorgesehen 'ID', 'Titel', 'Kursleiter', 'Beginn'. Durch Auswahl einer entsprechenden Checkbox ist es möglich die Suche auf kostenlose Kurse zu beschränken.
				\item \Func{Für Kursbenachrichtigungen eintragen}
					Der registrierte Nutzer hat die Möglichkeit sich bei der Anmeldung zu einem Kurs(\FuncBlue{/F20-120/}) durch Auswahl einer Checkbox in der Detailansicht des Kurses für Benachrichtigungen zu diesem Kurs zu registrieren oder nicht. Per Default ist die Checkbox nicht ausgewählt.
			\end{itemize}

		\subsection{Bezahlfunktionen}
			\begin{itemize}
				\item \Func{Guthabenkonto aufladen - online} \hypertarget{kontoAufladenOn}
					Der registrierte Nutzer besitzt ein Guthabenkonto. Dieses kann er online mittels Kreditkarte aufladen. Dafür betätigt der Benutzer die Schaltfläche 'Konto aufladen' auf seiner Profilseite. Er wird auf die 'Konto aufladen' - Seite weitergeleitet. Dort kann er seine Kreditkartendaten und den aufzuladenden Betrag eingeben. Die Kreditkartenabwicklung erfolgt dabei über die Infosun-Bank. Die Kreditkartendaten werden mittels SSL übertragen. Der aufgeladene Betrag ist sofort auf dem Guthabenkonto verfügbar und kann auch sofort verwendet werden. Aus Sicherheitsgründen werden keine Kreditkartendaten im System gespeichert.
				\item \Func{Guthabenkonto aufladen - offline} \hypertarget{kontoAufladenOff}
					Der registrierte Nutzer kann sein Guthabenkonto auch offline aufladen. Dies ist mittels Überweisung oder Barzahlung an den Betreiber oder einen Kursleiter möglich. 
					Der Nutzer überweist dafür den Betrag unter Angabe seiner Identifikationsnummer und seines Namens  auf das Konto des Betreibers oder bezahlt bar. Der bezahlte Betrag ist erst auf dem Guthabenkonto des registrierten Benutzers vorhanden, wenn er vom Betreiber  gutgeschrieben wurde(vgl. \FuncBlue{/F40-130/}).
				\item \Func{Kurseinheiten bezahlen}
					Die Bezahlung von Kurseinheiten wird automatisch bei der Anmeldung zur Kurseinheit erledigt. Der registrierte Nutzer kann sich also auch nur zu der Kurseinheit anmelden, wenn sein Guthaben auf dem Konto für den Preis der Einheit ausreicht bzw. sein Kontostand nicht den vom Betreiber festgelegten Überziehungskredit überschreitet(vgl. \FuncBlue{/F40-140/}).
				\item \Func{Rückbuchung des Kurspreises bei Abmeldung oder Kursausfall}
					Meldet sich der registrierte Nutzer von einer Kurseinheit ab oder fällt die Kurseinheit aufgrund zu geringer Teilnehmerzahl oder wegen eines anderen Grunds aus, so wird der Kaufpreis der Einheit automatisch auf das Konto des registrierten Nutzers zurückgebucht.
			\end{itemize}	
			
		\subsection{Terminplaner}
			\begin{itemize}
				\item \FuncW{Terminplaner.} \hypertarget{terminplaner}
					Der registrierte Benutzer hat einen persönlichen Terminplaner. Dieser wird in der Wochenansicht dargestellt. Der Terminplaner unterstützt nur die Anzeige von einstündigen Slots. Ist ein Termin kürzer, wird dennoch der Slot für die ganze Stunde belegt. Ist der Termin länger als eine Stunde werden dementsprechend mehrere einstündige Slots belegt.
					Außerdem ist darauf zu achten, dass die Slots jeweils zur vollen Stunde beginnen und zur nächsten vollen Stunde enden. (Analog dem Prinzip eines Stundenplans)
				\item \FuncW{Anzeigen des Terminplaners.}
					Der persönliche Terminplaner des registrierten Benutzers kann angezeigt werden. 
				\item \FuncW{Automatische Eintragung von angemeldeten Kurseinheiten.}
					Die angemeldeten Kurseinheiten des registrierten Benutzers werden automatisch in dessen Terminplaner eingetragen.
				\item \FuncW{Automatische Entfernung von abgemeldeten Kurseinheiten.}	
					Meldet sich der registrierte Benutzer aus einer Kurseinheit ab oder fällt eine Kurseinheit aus, so wird diese aus dem Terminplaner gelöscht.
				\end{itemize}   

	\section{Kursleiter}
	Jedem Kursleiter stehen auch alle Funktionen eines registrierten Benutzers zur Verfügung. Darüber hinaus kann ein Kursleiter auf folgende Funktionen zurückgreifen.
	\stepHc
		\subsection{Kurse}
			\begin{itemize}
				\item \Func{Eigene Kurse anzeigen.}
					Auf der Kursleiterseite werden die Kurse, welche der Kursleiter leitet, gesammelt angezeigt.
				\item \Func{Eigenen Kurs editieren.} \hypertarget{kursEdit}
					Der Kursleiter hat die Möglichkeit, die Kurse, die er  leitet, zu editieren. Dafür klickt er in seinen Trainingskursen den gewünschten Kurs an und wird auf die zugehörige 'Kursdetails' - Seite weitergeleitet. Durch Betätigen der 'Bearbeiten' - Schaltfläche kann er die Kursdaten bearbeiten. Er kann alle angezeigten Kursdaten verändern mit Ausnahme der vom System generierten Identifikationsnummer. Durch Betätigen der 
				\item \Func{Kurseinheit anlegen.} \hypertarget{kurseinheitAnlegen}
					Der Kursleiter hat die Möglichkeit Kurseinheiten für einen Kurs anzulegen, welchen er leitet. Dafür wechselt er in die 'Kursdetails' - Oberfläche. Dort betätigt er die Schaltfläche 'Kurseinheit anlegen'. Er wird auf eine Seite weitergeleitet, auf welcher er die Daten der Kurseinheit(\FuncBlue{/D30-10/}) eintragen kann. Das System generiert für die Kurseinheit eine im System eindeutige Identifikationsnummer. Durch Betätigen der 'Speichern' - Schaltfläche wird die Kurseinheit angelegt. Die Kurseinheiten können nur in dem Zeitraum zwischen Startdatum und Enddatum des zugehörigen Kurses stattfinden.
				\item \Func{Kurseinheit editieren.} \hypertarget{kurseinheitEdit}
					Der Kursleiter hat die Möglichkeit Kurseinheiten der eigenen Kurse zu editieren. Dafür klickt er in der Detailansicht des Kurses die 'Bearbeiten' - Schaltfläche der Kurseinheit an und wird auf die 'Kurseinheit bearbeiten' - Oberfläche weitergeleitet. Hier kann der Kursleiter die Daten der Kurseinheit(\FuncBlue{/D30-10/}) mit Ausnahme der vom System generierten Identifikationsnummer bearbeiten.
				\item \Func{Teilnehmer zu Kurseinheit hinzufügen.}
					Der Kursleiter hat die Möglichkeit durch Eingabe der Identifikationsnummer, des Namens und des Vornamens eines Benutzers diesen als Teilnehmer zu einer einzelnen Kurseinheit einer seiner Kurse einzutragen. Dafür muss der Kursleiter auf die 'Kurseinheit bearbeiten' - Seite der entsprechenden Kurseinheit wechseln.
				\item \Func{Teilnehmer aus Kurseinheit entfernen.}
				    Der Kursleiter hat die Möglichkeit Teilnehmer aus einer Kurseinheit eines von ihm geleiteten Kurses zu entfernen. Dafür muss der Kursleiter auf die 'Kurseinheit bearbeiten' - Seite der entsprechenden Kurseinheit wechseln. Dort wählt er den zu löschenden Teilnehmer aus der Teilnehmerliste der Kurseinheit aus und betätigt die 'Löschen' - Schaltfläche. Es besteht auch die Möglichkeit mehrere Teilnehmer gleichzeitig zu löschen.
				\item \Func{Kurseinheit löschen.} \hypertarget{kurseinheitLoeschen}
					Der Kursleiter kann Kurseinheiten eigener Kurse löschen. Dafür wechselt er durch Betätigen der 'Bearbeiten' - Schaltfläche der zu löschenden Kurseinheit auf die 'Kurseinheit bearbeiten' - Oberfläche. Durch Betätigen der 'Löschen' - Schaltfläche kann er die Kurseinheit löschen(vgl. \FuncBlue{/L10-50/}).
				\item \Func{Teilnehmer aus Kurs entfernen.} \hypertarget{kursUserDel}
					Der Kursleiter kann einzelne Teilnehmer aus einem seiner eigenen Kurse entfernen. Dafür wechselt er über die 'Kursdetails' - Oberfläche auf die Teilnehmerliste des Kurses. Dort wählt er den zu entfernenden Teilnehmer aus und entfernt diesen durch Betätigen der 'Löschen' - Schaltfläche aus dem Kurs(vgl. \FuncBlue{/L10-50/}). Es besteht auch die Möglichkeit mehrere Teilnehmer gleichzeitig zu entfernen.
			\end{itemize}

		\subsection{Benachrichtigungen}
			\begin{itemize}
				\item \Func{Benachrichtigung an Benutzergruppe senden.} \hypertarget{kursUserMsg}
					Der Kursleiter kann Benachrichtigungen an alle Teilnehmer seines Kurses oder nur an die Teilnehmer einer Kurseinheit seines Kurses senden. Voraussetzung dafür, dass die Teilnehmer Benachrichtigungen erhalten, ist deren Registrierung für die Kursbenachrichtigungen bei der Anmeldung zum Kurs.
				\item \Func{Benachrichtigung bei Kursänderung senden.}	
					Nimmt der Kursleiter Änderungen an einem seiner Kurse oder seiner Kurseinheiten vor, so kann er entscheiden, ob er alle Kursteilnehmer, nur die Teilnehmer der betreffenden Kurseinheit oder keinen Teilnehmer benachrichtigt. Voraussetzung dafür, dass die Teilnehmer Benachrichtigungen erhalten, ist deren Registrierung für die Kursbenachrichtigungen bei der Anmeldung zum Kurs.
			\end{itemize}

		\subsection{Verwaltung}
			\begin{itemize}
				\item \Func{Account aktivieren} \hypertarget{addUserKL}
					Der Kursleiter hat die Möglichkeit Accounts zu aktivieren, falls dies durch den Administrator in den Einstellungen für die Accountaktivierung so festgelegt wurde(vgl. \FuncBlue{/F40-160/}). Ist dies der Fall, so wird der Kursleiter durch Betätigen der 'Benutzer aktivieren' - Schaltfläche auf die 'Benutzer aktivieren' - Seite weitergeleitet. Dort findet er eine Auflistung aller Benutzer, die ihre E-Mailadresse bereits verifiziert haben, aber deren Account noch nicht aktiviert wurde. Der Kursleiter kann nun den Benutzer auswählen, dessen Account er aktivieren will. Durch Betätigen der 'Benutzer aktivieren' - Schaltfläche wird der Account des Benutzers aktiviert und dieser kann sich nun im System einloggen. 
					Es besteht auch die Möglichkeit mehrere Accounts gleichzeitig zu aktivieren. Nach Aktivierung der Benutzeraccounts erhalten die entsprechenden Nutzer eine automatische E-Mail, dass sie sich ab sofort im System anmelden können.
			\end{itemize}

	\section{Systemadministratoren}
	Jedem Systemadministrator stehen auch alle Funktionen eines Kursleiters zur Verfügung. Darüber hinaus kann ein Systemadministrator auf folgende Funktionen zurückgreifen.
	\stepHc
		\subsection{Kurse}
			\begin{itemize}
				\item \Func{Kurs anlegen.} \hypertarget{kursErstellen}
					Der Systemadministrator kann einen neuen Kurs anlegen. Dafür betätigt der Administrator die entsprechende Schaltfläche auf der Administratorseite und wird auf die 'Neuen Kurs anlegen' - Seite weitergeleitet. Die Anwendung generiert automatisch eine für den Kurs im System eindeutige Identifikationsnummer. Der Administrator hat die Möglichkeit die Daten des Kurses(\FuncBlue{/D20-10/}), wie etwa eine Beschreibung, einzugeben. Notwendige Eingaben sind der Startdatum, der Enddatum und mindestens ein Kursleiter für den Kurs. Innerhalb des Zeitraums von Startdatum bis Enddatum können Kurseinheiten stattfinden.
				\item \Func{Kurse verwalten.} \hypertarget{kursVerwalt}
					Der Systemadministrator kann sich alle Kurse anzeigen lassen. Dafür betätigt er die Schaltfläche 'Kurse verwalten' und wird auf die 'Kursangebot durchsuchen' - Oberfläche weitergeleitet. Dort kann er sich das gesamte Kursangebot anzeigen lassen oder nach bestimmten Kursen bezüglich verschiedener Suchparameter suchen. Durch Betätigen der Schaltfläche 'Suchen' werden die gesuchten Kurse angezeigt. Durch Klicken auf einen Kurs wird der Administrator auf die zugehörige 'Kursdetails' - Seite weitergeleitet.
				\item \Func{Kurs bearbeiten.} \hypertarget{kursBearbeiten}
					Der Systemadministrator hat die Möglichkeit einen beliebigen Kurs zu bearbeiten. Dafür wechselt er in die Detailansicht des zu editierenden Kurses. Durch Betätigen der 'Bearbeiten' - Schaltfläche  kann der Administrator nun die angezeigten Kursdaten editieren(vgl. \FuncBlue{/D20-10/}). Einzige Ausnahme ist die vom System generierte Identifikationsnummer des Kurses. Des Weiteren kann er Kurseinheiten anlegen, löschen oder editieren, Kursleiter zum Kurs hinzufügen oder entfernen und die Teilnehmer des Kurses editieren.
				\item \Func{Kurs löschen.} \hypertarget{kursLoeschen}
					Der Systemadministrator hat die Möglichkeit einen beliebigen Kurs zu löschen. Dafür wechselt er auf die 'Kursdetails' - Seite des zu löschenden Kurses. Durch Betätigen der 'Kurs löschen' - Schaltfläche wird der Kurs und damit auch alle zugehörigen Kurseinheiten gelöscht(vgl. \FuncBlue{/L10-50/}).
			\end{itemize}

	\subsection{Benutzer}
			\begin{itemize}
				\item \Func{Neuen Benutzer anlegen.} \hypertarget{nutzerAnlegen}
					Der Systemadministrator hat die Möglichkeit einen neuen Benutzer anzulegen. Dafür betätigt er die Schaltfläche 'Benutzer anlegen' auf der Administratorseite und wird auf die Seite 'Benutzer anlegen' weitergeleitet. 
					Die Anwendung generiert automatisch eine im System eindeutige Identifikationsnummer für den Benutzer. Der Systemadministrator kann die Benutzerdaten eingeben, ein Initialpasswort setzen und die Rolle des Benutzers festlegen. Die Benutzerrolle wird mittels einer Drop-Down - Liste aus den Möglichkeiten 'Kursleiter', 'Administrator','Benutzer' ausgewählt. Durch Betätigen der Schaltfläche 'Speichern' wird der neue Nutzer angelegt. Der vom Administrator erstellte Benutzer kann sich sofort ins System einloggen und muss nicht mehr aktiviert werden.
				\item \Func{Benutzer verwalten.} \hypertarget{nutzerVerwalt}
					Der Systemadministrator besitzt eine Übersicht über alle Benutzer. Diese erreicht er über die Schaltfläche 'Benutzer suchen' der Administratorseite. Auf der 'Benutzer durchsuchen' - Seite hat der Administrator die Möglichkeit sich alle Benutzer anzeigen zu lassen oder nach Benutzern zu suchen.  Dafür wählt er aus einer Drop - Down - Liste einen Parameter aus und gibt in ein Eingabefeld seinen Suchbegriff ein. Durch Betätigen der Schaltfläche 'Suchen' werden die gesuchten Benutzer angezeigt. Durch Klicken auf einen Benutzer wird der Administrator auf die 'Profil' - Seite des Benutzers weitergeleitet.
				\item \Func{Benutzer bearbeiten.} \hypertarget{nutzerBearbeiten}
					Der Systemadministrator hat die Möglichkeit einen beliebigen registrierten Nutzer zu bearbeiten. Dafür wechselt er auf die 'Profil' - Seite des zu bearbeitenden Nutzers. Durch Betätigen der 'Bearbeiten' - Schaltfläche kann er Administrator nun die Benutzerdaten editieren oder die Benutzerrolle verändern. Der Administrator kann alle Benutzerdaten verändern mit Ausnahme der vom System generierten Identifikationsnummer des Benutzers. Durch Betätigen der 'Speichern' - Schaltfläche werden die Daten gespeichert.
				\item \Func{Benutzer löschen.} \hypertarget{nutzerLoeschen}
					Der Systemadministrator hat die Möglichkeit einen beliebigen Benutzer zu löschen. Dafür wechselt er auf die 'Profil' - Seite des zu löschenden Benutzers. Durch Betätigen der 'Benutzer löschen' - Schaltfläche wird der Benutzer gelöscht(vgl. \FuncBlue{/L10-50/}). Nur der Administrator hat die Möglichkeit Benutzer zu löschen. Andernfalls könnte ein \glqq Nicht - Administrator\grqq \ einen Benutzer und damit auch sein Guthaben bzw. seine eventuellen Schulden löschen.
				\item \Func{Benutzerrolle verändern.} \hypertarget{nutzerRolle}
					Der Systemadministrator hat die Möglichkeit die Rolle jedes Nutzers zu verändern. Als Benutzerrollen sind 'Benutzer', 'Kursleiter' und 'Administrator' vorgesehen. Die Benutzerrolle kann auf der 'Benutzer bearbeiten' - Seite verändert werden.
				\item \Func{Benutzer suchen.} \hypertarget{nutzerSuchen}
					Der Systemadministrator hat die Möglichkeit nach Benutzern zu suchen. Dafür wählt er auf er Administratorseite die entsprechende Schaltfläche aus. Er wird auf die 'Benutzer durchsuchen' - Seite weitergeleitet, wo er die entsprechenden Suchparameter auswählen und seinen Suchbegriff eingeben kann. Durch Betätigen der Schaltfläche 'Suchen' wird das Suchergebnis angezeigt.
				\item \Func{Account aktivieren.}
				    Der Administrator kann genauso wie der Kursleiter Accounts aktivieren(vgl.\FuncBlue{/F30-120/}), falls dies vom Administrator so eingestellt wurde. Der Administrator wählt dafür die 'Benutzer verwalten' - Schaltfläche auf der Administratorseite aus und wird auf die 'Benutzer durchsuchen' - Seite weitergeleitet. Dort wählt er aus der Drop-Down-Liste die Option 'Nicht aktiviert' aus. Die nicht aktivierten Benutzer werden nun angezeigt. Zusätzlich werden nun auf der Seite die entsprechenden Schaltflächen angezeigt um die Benutzeraccounts zu aktivieren.
			\end{itemize}

		\subsection{Verwaltung}
			\begin{itemize}
				\item \Func{Statistiken anzeigen} \hypertarget{statistik}
					Der Systemadministrator hat die Möglichkeit sich Statistiken bezüglich der Einnahmen anzusehen. Er kann dafür aus 	verschiedenen Parametern wählen, welche Statistik angezeigt werden soll. Die Auswahl der Parameter erfolgt über eine Drop-Down - Liste. Vorgesehene Parameter für die Statistik sind Einnahmen pro 'Tag', pro 'Woche', pro 'Kursleiter', pro 'Kurs'. Bei Wahl eines der Parameter 'Tag', 'Woche' oder 'Kursleiter' ist ein graphische Darstellung in Form eines Säulendiagramms vorgesehen. Bei Parameter 'Tag' werden die Einnahmen der letzten sieben Tage angezeigt, bei Parameter 'Woche' die Einnahmen der letzten vier Wochen. Bei Wahl des Parameters 'Kurs' ist aufgrund der Menge von Kursen eine tabellarische Anzeige vorgesehen. Durch Betätigen der Schaltfläche 'Anzeigen' wird die gewählte Statistik angezeigt.
				\item \Func{Guthabenkonto aufladen} \hypertarget{guthabenAuf}
					Der Systemadministrator hat die Möglichkeit das Guthabenkonto von  registrierten Nutzern manuell aufzuladen. Damit kann er \glqq Offline - Aufladungen\grqq(\FuncBlue{/F20-190/})  des Guthabenkontos eines registrierten Nutzers auf dessen Konto im System weitergeben. Für die Kontoaufladung gibt der Administrator auf der Administratorseite die Identifikationsnummer des Benutzers, dessen Namen und den aufzuladenden Betrag ein. Durch Betätigen der Schaltfläche 'Konto aufladen' wird das Konto des Benutzers mit der eingegebenen Identifikationsnummer um den eingegebenen Betrag aufgeladen. Ab diesem Zeitpunkt kann der registrierte Benutzer das Guthaben verwenden. Das Aufladen des Guthabenkontos ist nur dann möglich, wenn der eingegebene Name auch zur Identifikationsnummer passt. So wird sichergestellt, dass durch Tippfehler bei der Identifikationsnummer nicht der falsche Nutzer das Guthaben erhält.
				\item \Func{Überziehungskredit festlegen.}
					Der Systemadministrator kann festlegen, ob und gegebenenfalls um welchen Betrag ein registrierter Benutzer sein Guthabenkonto überziehen darf. Der Überziehungskredit kann nur global für alle Nutzer festgelegt werden. Dafür trägt der Systemadministrator auf der Administratorseite den Betrag, der maximal überzogen werden darf in das dafür vorgesehene Eingabefeld ein und bestätigt die Eingabe durch Betätigen der 'Speichern'-Schaltfläche.
					Falls kein Überziehungskredit gewährt werden soll, wird als Betrag einfach 0,00 eingegeben. Dieser Wert ist auch per Default gesetzt.
			\end{itemize}
			
		\subsection{Systemanpassung}
			\begin{itemize}
				\item \Func{Logo hochladen.} \hypertarget{logo}
					Das Logo kann vom Systemadministrator editiert werden. Durch Klicken der 'Durchsuchen' - Schaltfläche auf der Administratorseite gibt der Administrator den Pfad einer .jpg - Datei an, welche durch Betätigen der 'Speichern' - Schaltfläche hochgeladen wird. Anforderungen an die .jpg - Datei bezüglich Größe und Auflösung werden während der Produktentwicklung festgelegt.
				\item \Func{Art der Accountaktivierung festlegen.} \hypertarget{regMod}
					Die Art der Aktivierung eines Accounts eines neuen Benutzers kann vom Systemadministrator festgelegt werden. Minimale Voraussetzung für die Aktivierung eines neuen Benutzeraccounts ist die Verifizierung der E-Mailadresse durch den Benutzer. Dafür betätigt der neue Benutzer den Verifizierungslink in der an die eingegebene E-Mailadresse gesandte Bestätigungsmail(vgl. \FuncBlue{/F10-20/}). Des Weiteren kann der Administrator festlegen, dass neue Benutzeraccounts zusätzlich manuell durch einen Kursleiter und/oder einen Administrator aktiviert werden müssen. Dafür wählt der Administrator auf der Administratorseite die entsprechende Option aus der Drop-Down-Liste und betätigt die 'Speichern' - Schaltfläche. Vorgesehene Auswahloptionen sind 'Nur E-Mailverifikation', 'E-Mailverifikation und Aktivierung durch Administrator', 'E-Mailverifikation und Aktivierung durch Kursleiter', 'E-Mailverifikation und Aktivierung durch Administrator oder Kursleiter'.
				\item \Func{Oberfläche editieren.}
					Die Oberfläche kann vom Systemadministrator editiert werden. Durch Klicken der 'Durchsuchen' - Schaltfläche gibt der Administrator den Pfad einer \gls{CSS} - Datei an, welche durch Betätigen der 'Speichern' - Schaltfläche hochgeladen wird. Die ausgewählte Datei muss den Dateinamen 'customStyle.css' haben. In dieser Datei kann zu bestimmten Tags ein neuer Stil definiert werden. 
					Die genauen Tags werden während der Produktentwicklung festgelegt. Der Administrator kann also die Oberfläche bezüglich bestimmter Punkte anpassen.  
					Falls keine 'customStyle.css' - Datei hochgeladen wird oder einige Tags in dieser Datei nicht definiert werden, wird ein im System definierter Default - Stil verwendet.
				\item \Func{Impressum editieren.} \hypertarget{impressumBea}
					Das Impressum kann vom Systemadministrator editiert werden. Durch Klicken der 'Impressum bearbeiten' - Schaltfläche wird der Administrator auf die 'Impressum bearbeiten' - Oberfläche weitergeleitet. Dort wird das Impressum in einem Eingabefeld angezeigt. Der Administrator kann nun Änderungen vornehmen. Durch Betätigen der 'Speichern' - Schaltfläche wird das geänderte Impressum gespeichert.
			\end{itemize}
		
\resetAllCounter
\renewcommand{\Func}[1]{\stepcounter{Lc}\textcolor{Blue}{\textbf{/D\arabic{Hc}0-\arabic{Lc}0/} #1} \\}	

\chapter{Produktdaten}
\label{Produktdaten}
Im folgenden wird beschrieben, welche Daten persistent gespeichert werden. Jeder Punkt stellt im Prinzip einen Datensatz dar.         

\section{System}
\Func {} Systemrelevante Informationen
\begin{itemize}
	\item Titel 
	\item CSS-Datei 
	\item Impressum 
	\item Logo
	\item AGB 
	\item Hilfe
	\item Überziehungskredit
	\item Art der Registrierung
	\item Konfigurationsdatei für Datenbankverbindung 
\end{itemize}


\subsection{SMTP-Server}
\Func {} Anmelde Daten für einen E-Maildienst von dem die Benachrichtigungen für Kursanmeldungen, Registrierung, Kursänderungen usw. verschickt werden sollen. 
\begin{itemize}
	\item Benutzername 
	\item Passwort 
	\item Authentifizierung 
	\item TLS 
	\item Host 
	\item Port 
\end{itemize}


\stepHc		    	
\section{Kurs}
\Func {} Alle Informationen über einen angebotenen Kurs. 
\begin{itemize}
	\item ID  (eindeutig)
	\item Daten
	\begin{itemize}
		\item Information zum Kurs
		\begin{itemize}
			\item Titel 
			\item Beschreibung 
			\item nächster Termin 
			\item Startdatum 
			\item Enddatum 
			\item Maximale Teilnehmerzahl
			\item Bild
			\item Kurseinheiten 		
		\end{itemize}
		\item Kontaktinformation
		\begin{itemize}
			\item Liste der Kursleiter
			\item Liste der Teilnehmer 
			\item Liste der zu benachrichtigenden Teilnehmer
		\end{itemize}	
	\end{itemize}
	
	
	
\end{itemize}
\stepHc		    
\section{Kurseinheit}
\Func {} Alle Informationen über eine angebotene Kurseinheit
\begin{itemize}
	\item ID (eindeutig)
	\item Daten
	\begin{itemize}
		\item Information zur Kurseinheit
		\begin{itemize}
			\item Titel 
			\item Beschreibung 
			\item Beginn 
			\item Ende 
			\item Ort 
			\item Preis 
			\item Maximale Teilnehmerzahl 
			\item Minimale Teilnehmerzahl 			
		\end{itemize}
		\item Kontaktinformation
		\begin{itemize}
			\item Kursleiter 
			\item Liste der Teilnehmer 
		\end{itemize}
	\end{itemize}
	
	
	
\end{itemize}   


\stepHc	   		
\section{Registrierte Benutzer}
\Func {} Alle Informationen über einen registrierten Benutzer
\begin{itemize}
	\item ID  (eindeutig)
	\item Kennung
	\begin{itemize}
		\item Benutzername (eindeutig)
		\item Passwort (gehasht)
		\item Salt (eindeutig)
	\end{itemize}
	\item Persönliche Daten
	\begin{itemize}
		\item Information zur eigenen Person	
		\begin{itemize}
			\item Vorname 
			\item Nachname 
			\item Geschlecht 
			\item Geburtsdatum 
			\item Bildverweis 
			\item Guthaben
		\end{itemize}
		\item Kontaktinformationen
		\begin{itemize}
			\item E-Mailadresse (eindeutig, gültig)
			\item Ort 
			\item Straße 
			\item Postleitzahl
		\end{itemize}
		\item Sonstige Daten
		\begin{itemize}
			\item eingeschriebene Kurse
			\item eingeschriebene Kurseinheiten
			\item Registrierung über Email abgeschlossen
			\item Registrierung über Admin abgeschlossen
		\end{itemize}	
	\end{itemize}
	
\end{itemize}


\stepHc	 
\section{Benutzerrolle}
\Func {} Speichert registrierte Benutzer gemäß ihrer Rolle im System. (3 verfügbare Benutzerrollen: Administrator, Kursleiter und Benutzer) 
\begin{itemize}
	\item ID (eindeutig)
	\item Gruppenname 
	\item Liste von Mitgliedern
\end{itemize}


\resetAllCounter
\renewcommand{\Func}[1]{\stepcounter{Lc}\textcolor{Blue}{\textbf{/L\arabic{Hc}0-\arabic{Lc}0/} #1} \\}	    

\chapter{Produktleistungen}


\section{Produktleistungen auf Serverseite}
\begin{itemize}
	\item \Func{} Zur Speicherung der \gls{persistent} Daten wird eine PostgreSQL Datenbank verwendet.
	\item \Func{} E-Mailbenachrichtigung, falls von Benutzer erwünscht, bei wichtigen Ereignissen (z.B Registrierung, Buchung einer Kurseinheit bzw Kurses, Terminänderung).
	\item \Func{} Das System hat Wartungszeit von 2 Stunden in der Woche bzw. 8 Stunden im Monat.
	\item \Func{} Kurse werden bis ein halbes Jahr nach Ablauf des Enddatums für den Statistikdienst gespeichert.
	\item \Func{} Löschen von Datensätzen wird kaskadenartig auf der Datenbank ausgeführt um Fehler bei Verlinkungen von Datensätzen zu vermeiden.
	\item \Func{} Das System stellt einen automatisch generierten Salt zur Verfügung um eine noch größere Sicherheit zu gewährleisten.
	\item \Func{} Die Anwendung stellt ein eigenes Bezahlsystem zur Verfügung, in der jeder User von seinem Prepaid Konto abbuchen kann.
\end{itemize}

\stepHc		
\section{Produktleistungen auf Clientseite}
\begin{itemize}
	\item \Func{} Bei Fehleingabe werden bereits eingetragene Daten wieder angezeigt, außer Passwörter und die falschen Benutzereingaben (z.B bei Registrierung).
	\item \Func{} Fehlermeldungen werden im Browser ausgegeben.
	\item \Func{} Passwörter werden mit SHA-256 \gls{hashen}.
	\item \Func{} Übertragung und Einsicht aller persönlichen Informationen erfolgen über eine SSL Verbindung.
	\item \Func{} Die Administratoren sind in der Lage registrierte Benutzer zu bearbeiten und ihnen neue Funktionen zuweisen (Kursleiter oder Administratoren).
	\item \Func{} Die Administratoren haben über eine selbst geschriebene CSS Datei Einfluss, auf ausgewählte Tags, die Optische Darstellung der Software zu beeinflussen. 
	\item \Func{} Unautorisierte Benutzeraktionen (z.B anonymer User will sich für Kurs eintragen) wird nicht angezeigt bzw. ausgeführt, stattdessen wird auf die Loginseite verwiesen.
	\item \Func{} An einem Client-Browser können nicht mehrere autorisierte verschiedene Sessions zeitgleich geführt werden.
	\item \Func{} Parallele Nutzung wird unterstützt und nach dem First Come First Serve Prinzip gehandhabt.
	\item \Func{} Um Skalierbarkeit zu gewährleisten wird Pagination bei der Anzeige von Kursen, Benutzern und Kurseinheiten verwendet. 
	\item \Func{} Das System stellt Validatoren zur Verfügung die Benutzereingaben auf Format Fehler überprüfen und eine Fehlermeldung bei verstoß ausgeben. (z.B. In einem Datumsfeld wird 33.02.1977 eingegeben)
	\item \Func{} Das System stellt einen Passwort Validator zur Verfügung, der das gewählte Passwort des Benutzers auf Sicherheitsstandards überprüft (mindestens 8 Zeichen, mindestens ein Sonderzeichen, mindestens eine Ziffer, Groß und Kleinbuchstaben). Außerdem dürfen im Passwort keine Umlaute und kein ß vorkommen.
	\item \Func{} Das System stellt einen E-Mail Validator zur Verfügung, der die vom Benutzer angegebene E-Mailadresse auf Gültigkeit überprüft.
	\item \Func{} Das System stellt einen Geburtstags Validator zur Verfügung, der das vom Benutzer angegebene Geburtsdatum kontrolliert ob es nicht in der Zukunft liegt. 
	\item \Func{} Das System stellt einen Überweisungs Validator zur Verfügung, der bei einer offline Überweisung des Trainers/Admins die angegebene BenutzerID auf Existenz überprüft.
	\item \Func{} Das System stellt einen Benutzernamen Validator zur Verfügung, der den vom anonymen User gewählten Benutzernamen bei der Registrierung auf bereits vergeben prüft.   
	\item \Func{} Verfügbarkeit der Räumlichkeit und der Terminkollision wird nicht geprüft.
	\item \Func{} Die Sprache des Systems kann auf Englisch oder Deutsch gesetzt werden. Benutzereingaben werden nicht berücksichtigt und werden unverändert angezeigt.
	\item \Func{} Der HTML 5 Code wird durch den Validator http://validator.w3.org überprüft.
	\item \Func{} Es werden keine Frames verwendet.
	\item \Func{} Um Zugriffskonflikte beim Nutzen des Systems mit mehreren Fenstern oder Tabs zu vermeiden wird mit Invalidate Session gearbeitet. (z.B. Ein Benutzer ist im System angemeldet und arbeitet mit zwei Tabs. Er meldet sich nun in einem Fenster ab und versucht darauf hin in dem anderen Fenster sich für einen Kurs anzumelden.)  
\end{itemize}

		
		
 
\chapter{Benutzeroberfläche}
	\begin{tiny}
		PC
	\end{tiny}
   
    \section{Screenshots}
    
	    \subsection{}
    
	    \subsection{}
    
    \section{Navigation}
	    In den Rechtecken der folgenden Grafik sind die verschiedenen Content-Bereiche der Webanwendung enthalten.
		    \newline
		    \newline
			    \includegraphics[scale=0.16]{logo/Navigation.pdf}
		    \newline
    
		\begin{itemize}
			\item Die Bereiche sind  nach Farben aufgeteilt (entsprechend der Autorisierung des Nutzers), sodass der jeweilige Bereich
		\begin{itemize}
		    \item Rot: für alle Nutzer,
		    \item Grün: nur für anonyme Nutzer (bei 'Registrierung' auch für Administratoren),
		    \item Blau: nur für angemeldete Nutzer,
		    \item Violett: nur für Kursleiter und Administratoren,
		    \item Gelb: nur für Administratoren
	    \end{itemize}
	    erreichbar ist.
		\item Die Rechtecke der Bereiche mit gepunkteter Umrandung befinden sich im  Header bzw. Footer der Webseite und sind global erreichbar.
		\item Die Rechtecke der Bereiche mit gestrichelter Umrandung befinden sich in der Navigationsleiste des angemeldetetn Nutzers. Diese beinhaltet, je nach Autorisierung, folgende Links zu Content-Bereiche:
		
		\begin{itemize}
		\item Registrierter Nutzer: 'Suche', 'Meine Kurse', 'Profil', 'Terminplaner'
		\item Kursleiter: 'Benutzersuche' (zusätzlich zur Navigationsleiste des registrierten Nutzers)
		\item Administrator: 'Adminverwaltung' (zusätzlich zur Navigationsleiste des Kursleiters)
		\end{itemize}		
		
	    \item Gestrichelte Pfeile bedeuten ein automatisches Weiterleiten.
	    \item Durch Ausloggen wird der Nutzer automatisch zur Startseite navigiert.
	    
	    \item Nachfolgend eine Erklärung zur Nummerierung der Pfeile im Diagramm:
	    
	    \begin{itemize}
	    \item 0: Nach dem Registrieren erhält der Nutzer eine E-Mail zur Verifikation. Klickt der Nutzer auf den darin enthaltenen Link wird er zur Startseite der Webanwendung navigiert.
	    \item 1: Durch Klicken auf den Link 'Registrierung' wird der Nutzer zu diesem Content-Bereich navigiert.
	    \item 2: Auf der Startseite befindet sich ein Link zum Bereich 'Suche', welcher neben einer Suchfunktion auch einen Auszug an angebotenen Kursen enthält.
	    \item 3: Durch Einloggen wird der Nutzer automatisch zum Bereich 'Meine Kurse' weitergeleitet (dies geschieht allgemein bei jedem besuchbaren Dialog für anonyme Nutzer).
	    \item 4: Nach der Suche nach Kursen in dem Content-Bereich 'Suche' erscheint eine Liste mit Ergebnissen. Der Nutzer kann anschließend entweder auf den gewünschten Kurs klicken und wird zum Bereich 'Kursdetails' des jeweiligen Kurses geleitet, oder auf den jeweiligen Kursleiter klicken und wird auf dessen Profil weitergeleitet. In diesen beiden Bereichen befindet sich ein Link, wodurch der Nutzer wieder zurück zur Ergebnisliste geleitet werden kann.
	    \item 5: Siehe Punkt 4.
	    \item 6: Im Bereich 'Meine Kurse' kann der Nutzer auf den gewünschten Kurs klicken und wird zum Bereich 'Kursdetails' des entsprechenden Kurses geleitet.
	    \item 7: Im Bereich 'Meine Kurse' kann der Nutzer auf den Kursleiter eines Kurses klicken und wird zu dessen Profil geleitet.
	    \item 8: Im Bereich 'Profil' befindet sich ein Link zum Dialog 'Kontoaufladung'
	    \item 9: Im Bereich 'Benutzersuche' bietet sich für Administratoren die Möglichkeit nach im System registrierten Nutzern zu suchen. Nach der Suche kann der Administrator in der Ergebnisliste auf den gewünschten Nutzer (entweder Kursleiter, Administrator oder normaler registrierter Nutzer) klicken und wird zu dessen Profil weitergeleitet.
	    \item 10: Siehe Punkt 9.
	    \item 11: Nach Klicken des Links 'Kurse verwalten' im Dialog 'Adminverwaltung' wird der Nutzer zum Bereich 'Suche' weitergeleitet. Hier kann der Nutzer nach den gewünschten Kursen suchen.
	    \item 12: Im Bereich 'Adminverwaltung' befindet sich ein Link zum Dialog 'Impressum bearbeiten'.
	    \item 13: Im Bereich 'Adminverwaltung' befindet sich ein Link zum Dialog 'Kurs anlegen'.
	    \item 14: Im Bereich 'Adminverwaltung' befindet sich ein Link zum Dialog 'Geldstatistik'.
	    \item 15: Im Bereich 'Adminverwaltung' befindet sich ein Link zum Dialog 'Benutzer anlegen'.
	    \item 16: Im Bereich 'Kursdetails' können Kursleiter sowie Administratoren zum Bearbeiten der Teilnehmer eines Kurses durch Klicken des entsprechenden Links zum Dialog 'Kursteilnehmer bearbeiten' weitergeleitet werden. Hier befindet sich ein Link, wodurch der Nutzer wieder zum Bereich 'Kursdetails' des Kurses zurückkehren kann.
	    \item 17: Im Bereich 'Kursdetails' führt ein Klick auf die entsprechende Kurseinheit zum Dialog 'Kurseinheit editieren'. Hier kann der Kursleiter oder Administrator die Teilnehmer sowie die Details der Kurseinheit bearbeiten. Durch einen Link kann der Nutzer wieder zur Ansicht 'Kursdetails' des Kurses zurückkehren.
	    \end{itemize}
	    
	    \end{itemize}
        
    \section{Bedienung}       
	    \begin{itemize}
			\item /B10/ Standardmäßig ist eine menüorientierte Bedienung vorgesehen
			\item /B20/ Die Bedienungsoberfläche ist auf Maus- und Tastaturbedienung ausgelegt
			\item /B30/ Fensterlayout, Dialogstruktur und Mausbedienung entsprechen dem Windows-Gestaltungs-Regelwerk
			\item /B40/ Sämtliche datenschutzrechtlich relevante Daten sind passwortgeschützt
			\item /B50/ Die Darstellung des Fensterinhalts passt sich der Fenstergröße an
		\end{itemize}
            
            
         
\chapter{Qualitätsbestimmungen}
	\begin{tiny}
		RS
	\end{tiny}\\
	Im Folgenden wird aufgeführt auf welche Qualitätsmerkmale der Webanwendung Wert gelegt wird. Diese orientieren sich an der Norm \textit{ISO/IEC~9126}, welche ein Modell zur Sicherstellung von Softwarequalität darstellt. In wie weit die verschiedenen Qualitätsmerkmale Einfluss auf die Webanwendung nehmen sollen, wird in der unten gezeigten Tabelle definiert.\\

	\begin{table}[h]
 
	    \begin{center}
		    \begin{tabular}{|l|c|}
			    \hline 
			    \rule[-1ex]{0pt}{2.5ex} \textbf{Qualitätskriterium} & \textbf{Stellenwert} \\ 
			    \hline 
			    \rule[-1ex]{0pt}{2.5ex} Funktionalität & sehr wichtig \\ 
			    \hline 
			    \rule[-1ex]{0pt}{2.5ex} Zuverlässigkeit & sehr wichtig \\ 
			    \hline 
			    \rule[-1ex]{0pt}{2.5ex} Benutzbarkeit & sehr wichtig \\ 
			    \hline 
			    \rule[-1ex]{0pt}{2.5ex} Effizienz & wichtig \\ 
			    \hline 
			    \rule[-1ex]{0pt}{2.5ex} Änderbarkeit & weniger wichtig \\ 
			    \hline 
			    \rule[-1ex]{0pt}{2.5ex} Übertragbarkeit & weniger wichtig \\ 
			    \hline   
		    \end{tabular}  
	    \end{center}
	    \caption{Qualitätsmerkmale und ihre zugewiesene Priorität}
	\end{table}

Der Stellenwert der Qualitätskriterien begründet sich wie folgt:
    	\begin{itemize}
    		\item Funktionalität: Alle genannten Produktfunktionen sollen implementiert werden und funktionieren, weshalb der Funktionalität eine sehr wichtige Bedeutung zukommt.
    		\item Zuverlässigkeit: Das System soll für den Anwender zu jeder Zeit verfügbar und nur mit geringen Wartungsaufwand verbunden sein, was eine hohe Zuverlässigkeit voraussetzt.
    		\item Benutzbarkeit: Die Anwendung soll von Nutzern einfach und intuitiv bedienbar sein ohne eine Hilfestellung zu benötigen, weshalb die Benutzbarkeit als sehr wichtig einzustufen ist.
    		\item Effizienz: Systemressourcen sind mit hohen Kosten verbunden, weshalb eine effiziente Anwendung programmiert werden soll. Andere Kriterien, wie beispielsweise Funktionalität, erhalten jedoch eine höhere Priorität.
    		\item Änderbarkeit: Um die Anwendung nach Fertigstellung erweitern zu können, ist ein gewisses Maß an Änderbarkeit notwendig. Dieser Punkt wird bei der Systementwicklung berücksichtigt, jedoch haben Kriterien, wie beispielsweise Zuverlässigkeit, einen höheren Stellenwert.
    		\item Übertragbarkeit: Die entwickelte Webanwendung soll in erster Linie auf den von uns benutzten und vergleichbaren Systemen laufen. Um beispielsweise ein anderes Datenbanksystem zu unterstützen ist Übertragbarkeit wichtig, jedoch wird bei der Entwicklung den anderen Kriterien ein deutlich höherer Stellenwert zugeschrieben.
    	\end{itemize}
	Darüber hinaus sollten noch folgende Qualitätsmerkmale genannt werden, welchen ein besonders hoher oder ein, zur Abgrenzung, geringer Stellenwert zugesprochen wird:
		\begin{itemize}
		 	\item Die Benutzung soll komplett im Browser möglich sein, wobei jedoch die Kompatibilität mit älteren Browsern als definiert nicht berücksichtigt wird.
		 	\item Personenbezogene Daten und im Speziellen die Bankverbindung sind vor unberechtigtem Zugriff von außen zu schützen.
		 	\item Geldtransaktionen sind ebenfalls besonders zu schützen, vor allem die \gls{Konsistenz} aller Bezahlaktionen.
 \end{itemize}
 
\resetAllCounter
\newcommand{\Test}[1]{\stepcounter{Lc}\textcolor{Brown}{\textbf{/T\arabic{Hc}0-\arabic{Lc}0/} #1} \\}
\newcommand{\RefFuncBlue}[1]{\textcolor{Blue}{\textbf{#1}}}
\newcommand{\RefFuncGreen}[1]{\textcolor{Green}{\textbf{#1}}}
\newcommand{\RefFuncBrown}[1]{\textcolor{Brown}{\textbf{#1}}}
\chapter{Globale Testszenarien und Testfälle}
 

	\section{Testfälle für den Systemadministrator ohne bestehenden Datensatz}
		\subsection{Setup}
			\begin{itemize}
				\item \Test{Anmeldung Administrator} 
			     Der Administrator ruft die Startseite auf und loggt sich hier über die Schaltfläche 'Anmelden' in dem sich öffnenden Drop-Down Menü mit dem vorgefertigten Benutzername 'admin1' und dem Passwort 'ersterAdmin1' ein und wird nach dem Drücken des sich im Drop-Down Menü befindenden Buttons 'Anmelden' auf die Seite 'Meine Kurse' weitergeleitet. 
				 
				\item \Test{Impressum bearbeiten}
				 Der Systemadministrator befindet sich auf der Seite 'Adminverwaltung', welche er über die Navigationsleiste ausgewählt hat, und klickt dort auf die Schaltfläche 'Impressum bearbeiten' und wird auf diese weitergeleitet. Auf der Seite 'Impressum bearbeiten' fügt er seinen Impressumstext in das Eingabefeld ein, in welchem als Default bereits ein vorgefertigtes Impressum der Firma ofCourse angegeben ist. Anschließend bestätigt der Admin die Eingabe durch Drücken des Button 'Speichern'. Mit dem Bestätigen der Änderungen am Impressum wird der Administrator wieder auch die Seite 'Adminverwaltung' geleitet. (vgl. \RefFuncBlue{/F40-180/})
				 
				
			\end{itemize}		
				
		\subsection{Erstellung und Verwaltung}
			\begin{itemize}
				\item \Test{Überziehungskredit festlegen} 
				 Auf der Seite 'Adminverwaltung' legt der Systemadministrator den Überziehungskredit auf 10 Euro fest, indem er im Textfeld den bereits genannten Betrag '10,00' eingibt und diesen durch Drücken der Schaltfläche 'Speichern' bestätigt. Vor dieser Änderung wurde im Textfeld der Default-Wert von '0,00' angezeigt, nun zeigt dieses Feld den Wert '10,00' an. Der Administrator befindet sich weiterhin auf der Seite 'Adminverwaltung'. (vgl. \RefFuncBlue{/F40-140/})
				
				\item \Test{Art der Accountaktivierung festlegen}
				 Der Administrator befindet sich auf der Seite 'Adminverwaltung'. Im Bereich 'Accountaktivierung' wird im Eingabefeld die voreingestellte minimale Voraussetzung 'Nur E-Mailverifikation' angezeigt. Dort wählt der Admin nach dem Klicken auf dieses Eingabefeld im sich öffnenden Drop-Down Menü die Aktivierungsart 'E-Mailverifikation und Aktivierung durch Administrator' aus und bestätigt seine Auswahl durch Drücken des Button 'Speichern'. Im Eingabefeld wird nun die ausgewählte Art 'E-Mailverifikation und Aktivierung durch Administrator' angezeigt. Der Administrator befindet sich weiterhin auf der Seite 'Adminverwaltung'. (vgl. \RefFuncBlue{/F40-160/})
				
				\item \Test{Neuen Benutzer anlegen} 
				 Auf der Seite 'Adminverwaltung' klickt der Administrator auf die Schaltfläche 'Benutzer anlegen' und wird auf die Seite 'Benutzer anlegen' weitergeleitet. Hier gibt der Admin folgende Daten ein:
					\begin{itemize}
						\item Benutzername: 'Basti3' 
						\item Passwort: 'poIuZt7R'
						\item Passwort: 'poIuZt7R'
						\item E-Mail: 'sebastian@nrschwarz.de'
						\item Anrede: 'Herr'
						\item Name: 'Schwarz'
						\item Vorname: 'Sebastian'
						\item Geburtstag: '04.12.1990'
						\item Straße: 'Musterstraße 1'
						\item Postleitzahl: '99999'
						\item Ort: 'Musterstadt'
					\end{itemize}
				 Als Rolle wird 'Kursleiter' festgelegt. Die BenutzerID wird automatisch generiert. Über den Button 'Durchsuchen' unter dem Dummy-Bild wählt der Admin ein Bild des Dateiformates .jpg aus seinem Verzeichnis aus. Dieses erscheint nun anstelle des Dummy-Bildes. Durch Klicken des Buttons 'Benutzer anlegen' bestätigt der Admin die Daten und der neue Benutzer ist angelegt. Der Administrator wird anschließend auf die Seite 'Adminverwaltung' weitergeleitet. (vgl. \RefFuncBlue{/F40-50/})
				
				\item \Test{Neuen Benutzer anlegen} 
				Auf der Seite 'Adminverwaltung' klickt der Administrator auf die Schaltfläche 'Benutzer anlegen' und wird auf die Seite 'Benutzer anlegen' weitergeleitet. Hier gibt der Admin folgende Daten ein:
				\begin{itemize}
					\item Benutzername: 'PatrickC' 
					\item Passwort: 'lKjHg789'
					\item Passwort: 'lKjHg789'
					\item E-Mail: 'strat-fan@web.de'
					\item Anrede: 'Herr'
					\item Name: 'Patrick'
					\item Vorname: 'Cretu'
					\item Geburtstag: '03.01.1989'
					\item Straße: 'Musterstraße 5'
					\item Postleitzahl: '99999'
					\item Ort: 'Musterstadt'
				\end{itemize}
				Als Rolle wird 'Benutzer' festgelegt. Die BenutzerID wird automatisch generiert. Durch Klicken des Buttons 'Benutzer anlegen' bestätigt der Admin die Daten und der neue Benutzer ist angelegt. Der Administrator wird anschließend auf die Seite 'Adminverwaltung' weitergeleitet. (vgl. \RefFuncBlue{/F40-50/})
				
				\item \Test{Neuen Benutzer anlegen} 
				 Auf der Seite 'Adminverwaltung' klickt der Administrator auf die Schaltfläche 'Benutzer anlegen' und wird auf die Seite 'Benutzer anlegen' weitergeleitet. 
					 \begin{itemize}
					 \item Hier gibt der Admin folgende Daten ein:
						\begin{itemize}
							\item Benutzername: 'Ricky1' 
							\item Passwort: 'vBn§'
							\item Passwort: 'vBn§'
							\item E-Mail: 'ricky.strohmeier@gmx.de'
							\item Anrede: 'Herr'
							\item Name: 'Strohmeier'
							\item Vorname: 'Ricky'
							\item Geburtstag: '40.10.1991'
							\item Straße: 'Musterstraße 2'
							\item Postleitzahl: '99999'
							\item Ort: 'Musterstadt'
						\end{itemize}
					 Als Rolle wird 'Benutzer' festgelegt und im Bereich Profilbild bleibt das Dummy-Bild bestehen. Da der 40.10.1991 nicht existiert wird nach dem Drücken der Schaltfläche 'Benutzer anlegen' eine Fehlermeldung ausgegeben, welche den Benutzer auf das falsche Datum hinweist und das entsprechende Eingabefeld wird geleert. Aufgrund des zu kurzen Passwortes und der Verwendung des unerlaubten Zeichens '§' wird der Benutzer des Weiteren mit einer Fehlermeldung darauf hingewiesen, dass sein Passwort nicht akzeptiert wird. Die beiden Passworteingabefelder werden ebenfalls geleert.
					 
					 \item Der Benutzer trägt nun im Eingabefeld des Geburtsdatums den 04.10.1991 ein. Außerdem gibt er ein anderes Passwort im ersten Passwortfeld ein: 'vBnMk89po'. Im zweiten Passwortfeld trägt er 'vBnMk89pO' ein. Nach dem Drücken der Schaltfläche 'Benutzer anlegen' wird der Benutzer mit einer Fehlermeldung darauf hingewiesen, dass seine Passwörter nicht übereinstimmen. Die beiden Passwortfelder werden wieder geleert.
					 
					 \item Der Benutzer gibt nun in beiden Passwortfeldern das Passwort 'vBnMk89po' ein. Durch Klicken des Buttons 'Benutzer anlegen' bestätigt der Admin die Daten und der neue Benutzer ist angelegt. Der Administrator wird anschließend auf die Seite 'Adminverwaltung' weitergeleitet.
					\end{itemize}
				(vgl. \RefFuncBlue{/F40-50/})
					
				\item \Test{Neuen Kurs anlegen}
			 	 Der Administrator befindet sich auf der Seite 'Adminverwaltung' und klickt auf die Schaltfläche 'Kurs anlegen'. Dadurch wird er auf die Seite 'Neuen Kurs anlegen' weitergeleitet.
				 	 \begin{itemize} 
					 	 \item Dort trägt er folgende Daten ein:
							\begin{itemize}
								\item Kursname: 'Yoga'
								\item Beschreibung: 'Bringe Körper, Geist und Seele in Einklang durch Yogaübungen. Jeder Teilnehmer sollte seine eigene Turnmatte mitnehmen.'
								\item Startzeitpunkt: '27.07.2015'
								\item Endzeitpunkt: '27.07.2016'
								\item minimale Teilnehmerzahl: '5'
							\end{itemize}
							Nach Drücken der Schaltfläche 'Kurs anlegen' wird der Administrator durch eine Fehlermeldung darauf hingewiesen, dass kein einziger Kursleiter eingetragen wurde.
							
						 \item Der Administrator trägt nun im Bereich 'Kursleiter hinzufügen' den Benutzernamen, die E-Mailadresse und die KursleiterID des Kursleiters 'Basti3' ein. (Daten vgl. \RefFuncBrown{/T10-50/}) Durch Klicken des Buttons 'Kurs anlegen' legt der Admin den neuen Kurs 'Yoga' an und wird anschließend auf die Seite 'Adminverwaltung' weitergeleitet. Die KursId wird automatisch vom System generiert. (vgl. \RefFuncBlue{/F40-10/})
					\end{itemize}
					
				\item \Test{Neuen Kurs anlegen}
				 Der Administrator befindet sich auf der Seite 'Adminverwaltung' und klickt auf die Schaltfläche 'Kurs anlegen'. Dadurch wird er auf die Seite 'Neuen Kurs anlegen' weitergeleitet.
					 \begin{itemize}
						\item Dort trägt er folgende Daten ein:
							\begin{itemize}
								\item Name: 'Standardtanz'
								\item Startzeitpunkt: '04.08.2015'
								\item Endzeitpunkt: '33.09.2015'
								\item minimale Teilnehmerzahl: '2'
							\end{itemize}
						Im Bereich 'Kursleiter hinzufügen' trägt der Administrator den Benutzernamen, die E-Mailadresse und die KursleiterID des Kursleiters 'Basti3' ein. (Daten vgl. \RefFuncBrown{/T10-50/}) Da der 33.09.2015 nicht existiert, wird der Administrator nach dem Drücken des Button 'Kurs anlegen' durch eine Fehlermeldung auf seine falsche Eingabe hingewiesen. Das Feld für den Endzeitpunkt wird geleert. 
						
						\item Der Administrator ändert den Endzeitpunkt in dem entsprechenden Eingabefeld auf den 30.09.2015 und fügt folgende Beschreibung in dem dafür vorgesehenen Eingabefeld hinzu: 'Grundkurs Standardtanz. Erlernt werden Tänze wie Walzer, Disco Fox, Cha-Cha-Cha, Slow Fox und Tango. Die Teilnahme als Paar ist wünschenswert.' Durch Klicken des Buttons 'Kurs anlegen' legt der Admin den neuen Kurs 'Standardtanz' an. Der Administrator wird anschließend auf die Seite 'Adminverwaltung' weitergeleitet. Die KursId wird automatisch vom System generiert. (vgl. \RefFuncBlue{/F40-10/}) 		
				 
					\end{itemize}
				 			
			\end{itemize}			
								
		\subsection{Systemdarstellung}
			\begin{itemize}
				 \item \Test{Logo hochladen} 
				 Auf der Seite 'Adminverwaltung' ändert der Systemadministrator das Logo, indem er im Bereich 'Logo ändern' auf 'Durchsuchen' klickt.
					\begin{itemize}
						\item Nach Auswahl der Musikdatei 'Test.mp3' drückt der Admin die Schaltfläche 'Speichern'. Durch eine Fehlermeldung wird der Administrator darauf hingwiesen, dass nur JPG-Dateien akzeptiert werden. Das vorherige Logo bleibt unverändert.
						
					 	\item Nach Auswahl der Bilddatei 'logo2.jpg' drückt der Admin die Schaltfläche 'Speichern', wodurch das Bild hochgeladen wird. Der Administrator befindet sich weiterhin auf der Seite 'Adminverwaltung'. Nach erneutem Laden dieser Seite wird das neue Logo auf der Oberfläche angezeigt.
					\end{itemize}
				 (vgl. \RefFuncBlue{/F40-150/})	
				
				 \item \Test{Oberfläche editieren}
				 Der Systemadministrator befindet sich auf der Seite 'Adminverwaltung' und klickt im Bereich 'Oberfläche anpassen' auf den Button 'Durchsuchen'. 
					\begin{itemize}
						 \item In dem sich öffnenden Fenster wählt der Admin die Datei mit dem Namen 'aenderung.css' aus und bestätigt durch Drücken der Schaltfläche 'Speichern'. Da dieser Dateiname nicht genehmigt ist wird bei erneutem Laden der Seite 'Adminverwaltung' der im System vordefinierte Default-Stil als Oberfläche angezeigt.
						 
						 \item Der Admin wählt erneut die Schaltfläche 'Durchsuchen' im Bereich 'Oberfläche anpassen' aus. In dem sich öffnenden Fenster wählt der Administrator nun die Datei mit dem Namen 'customStyle.css' aus. Nach dem Drücken des Buttons 'Speichern' wird diese Datei hochgeladen. Durch Drücken der Taste 'F5' wird die aktuelle Seite 'Adminverwaltung' neu geladen und die Änderungen durch die gewählte .css-Datei wird sichtbar. 
					\end{itemize}
				 (vgl. \RefFuncBlue{/F40-170/})
				 
			\end{itemize}
			
	\section{Testfälle für den Kursleiter}
		\stepHc
		\begin{itemize}
			\item \Test{Eigenen Kurs editieren} 
			Der Kursleiter 'Basti3' befindet sich auf der Seite 'Meine Kurse' und klickt in seinen Trainingskursen den Kurs 'Yoga' an. Dadurch wird er auf die 'Kursdetails'-Seite weitergeleitet, auf der bereits die durch den Systemadministrator angegebenen Daten eingetragen sind. (vgl. \RefFuncBrown{/T10-70/}) Der Kursleiter drückt auf den Button 'Bearbeiten', um die Kursdaten zu ändern:
			\begin{itemize}
				\item Der Button 'Bearbeiten' ändert sich nun in den Button 'Speichern' und der Kursleiter kann als neues Startdatum in das dafür vorgesehene Eingabefeld den '32.08.2015' eintragen. Nach dem Drücken der Schaltfläche 'Speichern' wird der Kursleiter durch eine Fehlermeldung darauf hingewiesen, dass das Datum falsch ist. Das Eingabefeld 'Startdatum' wird geleert.
				
				\item Der Kursleiter ändert den Startzeitpunkt des Kurses nun auf den '05.08.2015' und bestätigt seine Änderung durch Drücken des Button 'Speichern'. Das neue Datum wird nun im entsprechenden Feld angezeigt und der Button 'Speichern' ändert sich wieder zum Button 'Bearbeiten'. Der Kursleiter bleibt weiterhin auf der Seite 'Kursdetails'. (vgl. \RefFuncBlue{/F30-20/})
			\end{itemize}
			
			\item \Test{Kurseinheit anlegen} 
			Der Kursleiter 'Basti3' befindet sich auf der Seite 'Meine Kurse' und klickt in seinen Trainingskursen den Kurs 'Yoga' an. Dadurch wird er auf die 'Kursdetails'-Seite weitergeleitet auf der bereits die durch den Systemadministrator angegebenen Daten eingetragen sind. (vgl. \RefFuncBrown{/T10-80/}) Nach Drücken des Button 'Kurseinheit anlegen' wird er auf die Seite 'Kurseinheit bearbeiten' weitergeleitet. Nach Drücken des Button 'Bearbeiten', welcher sich dadurch zum Button 'Speichern' ändert, können die einzelnen Eingabefelder bearbeitet werden und der Kursleiter legt folgende Kurseinheit an:
			\begin{itemize}
				\item Termin: '05.08.2015'
				\item Beginn: 17:00
				\item Ende: 19:00
				\item Ort: 'Turnhalle'
				\item Beschreibung: 'Bringe Körper, Geist und Seele in Einklang durch Yogaübungen. Jeder Teilnehmer sollte seine eigene Turnmatte mitnehmen.'
				\item Preis: '00,00'
				\item Kursleiter: 'Basti3'
				\item Mindestteilnehmerzahl: '5'
				\item Maximale Teilnehmerzahl: '30'
				\item KursId: 'autogenerierte Id des Kurses'
			\end{itemize}	
			Anschließend bestätigt der Kursleiter die Kurseinheit durch Drücken des Buttons 'Speichern', welcher sich dadurch wieder zum Button 'Bearbeiten ändert. Der Kursleiter wird auf die Seite 'Kursdetails' weitergeleitet, auf der im Bereich 'Kurseinheiten' nun die eben angelegte Kurseinheit angezeigt wird. (vgl. \RefFuncBlue{/F30-40/})
			
			\item \Test{Kurseinheit anlegen} 
			Analog zum vorangehenden Testfall \RefFuncBrown{/T20-20/} werden zum selben Kurs 'Yoga' noch vier weitere Kurseinheiten zu den nachfolgenden Terminen angelegt:
			\begin{itemize}
				\item 12.08.2015
				\item 19.08.2015
				\item 26.08.2015
				\item 02.09.2015
			\end{itemize}		
			Alle anderen Daten werden analog zum Testfall \RefFuncBrown{/T20-20/} eingetragen.
			
			\item \Test{Kurseinheit anlegen} 
			Der Kursleiter 'Basti3' befindet sich auf der Seite 'Meine Kurse' und klickt in seinen Trainingskursen den Kurs 'Standardtanz' an. Dadurch wird er auf die 'Kursdetails'-Seite weitergeleitet auf der bereits die durch den Systemadministrator angegebenen Daten eingetragen sind. (vgl. \RefFuncBrown{/T10-90/}) Nach Drücken des Button 'Kurseinheit anlegen' wird er auf die Seite 'Kurseinheit bearbeiten' weitergeleitet. Nach Drücken des Button 'Bearbeiten', welcher sich dadurch zum Button 'Speichern' ändert, können die einzelnen Eingabefelder bearbeitet werden und der Kursleiter legt folgende Kurseinheit an:
			\begin{itemize}
				\item Termin: '03.08.2015'
				\item Beginn: 18:00
				\item Ende: 21:00
				\item Ort: 'Turnhalle'
				\item Beschreibung: 'Grundkurs Standardtanz. Erlernt werden Tänze wie Walzer, Disco Fox, Cha-Cha-Cha, Slow Fox und Tango. Die Teilnahme als Paar ist wünschenswert.'
				\item Preis: '15,00'
				\item Kursleiter: 'Basti3'
				\item Mindestteilnehmerzahl: '2'
				\item Maximale Teilnehmerzahl: '10'
				\item KursId: 'autogenerierte Id des Kurses'
			\end{itemize}	
			Anschließend bestätigt der Kursleiter die Kurseinheit durch Drücken des Buttons 'Speichern', welcher sich dadurch wieder zum Button 'Bearbeiten ändert. Der Kursleiter wird auf die Seite 'Kursdetails' weitergeleitet, auf der im Bereich 'Kurseinheiten' nun die eben angelegte Kurseinheit angezeigt wird. (vgl. \RefFuncBlue{/F30-40/})
			
			\item \Test{Kurseinheit anlegen} 
			Analog zum vorangehenden Testfall \RefFuncBrown{/T20-20/} werden zum selben Kurs 'Standardtanz' noch fünf weitere Kurseinheiten zu den nachfolgenden Terminen angelegt:
			\begin{itemize}
				\item 10.08.2015
				\item 17.08.2015
				\item 24.08.2015
				\item 31.08.2015
				\item 07.09.2015
			\end{itemize}		
			Alle anderen Daten werden analog zum Testfall \RefFuncBrown{/T20-40/} eingetragen.
			
			\item \Test{Kurseinheiten editieren}
			Der Kursleiter 'Basti3' befindet sich auf der Seite 'Meine Kurse' und klickt in seinen Trainingskursen den Kurs 'Standardtanz' an, wodurch er auf die 'Kursdetails'-Seite weitergeleitet wird, auf der bereits die durch den Systemadministrator angegebenen Daten (vgl. \RefFuncBrown{/T10-90/}) eingetragen sind. Des Weiteren werden die in den Testfällen \RefFuncBrown{/T20-40/} und \RefFuncBrown{/T20-50/} angelegten Kurseinheiten angezeigt. Der Kursleiter klickt auf den Button 'Bearbeiten', der sich neben der Kurseinheit mit dem Termin '03.08.2015' befindet. Dadurch wird er auf die Seite 'Kurseinheit bearbeiten' weitergeleitet, auf der bereits die im Testfall \RefFuncBrown{/T20-40/} angelegten Daten eingetragen sind. Nach Drücken des Button 'Bearbeiten', welcher sich dadurch zum Button 'Speichern' ändert, können die einzelnen Eingabefelder bearbeitet werden. 
			\begin{itemize}
				\item Der Kursleiter ändert den Termin auf den '01.10.2015' und drückt den Button 'Speichern'. Durch eine Fehlermeldung wird der Kursleiter darauf hingewiesen, dass sich der geänderte Termin außerhalb des Kurszeitraumes (vgl. \RefFuncBrown{/T10-80/}) befindet. Das Eingabefeld des Termins wird geleert. Der Button 'Speichern' ändert sich nicht.
				
				\item Nun ändert der Kursleiter den Termin auf den '05.13.2015' und drückt den Button 'Speichern'. Durch eine Fehlermeldung wird der Kursleiter darauf hingewiesen, dass der geänderte Termin kein gültiges Datum ist. Das Eingabefeld des Termins wird geleert. Der Button 'Speichern' ändert sich nicht.
				
				\item Der Kursleiter ändert den Termin auf den '04.08.2015' und die Maximale Teilnehmerzahl auf '2'. Anschließend drückt er den Button 'Speichern'. Da dieses Datum akzeptiert wird, wird 'Basti3' auf die Seite 'Kursdetails' des Kurses 'Standardtanz' weitergeleitet, auf der nun im Bereich 'Kurseinheiten' die entsprechende Einheit mit dem geänderten Termin angezeigt wird.
			\end{itemize}
			(vgl. \RefFuncBlue{/F30-50/})
			
			\item \Test{Kurseinheit löschen}
			Der Kursleiter 'Basti3' befindet sich auf der Seite 'Kursdetails' des Kurses 'Standardtanz'. Auf dieser befinden sich bereits die durch den Systemadministrator angegebenen Daten (vgl. \RefFuncBrown{/T10-90/}), sowie die angelegten Kurseinheiten aus dem Testfall \RefFuncBrown{/T20-60/}. Der Kursleiter wählt auf dieser Seite die Kurseinheit am '24.08.2015' durch Drücken des Button 'Bearbeiten' aus und wird auf die Seite 'Kurseinheit bearbeiten' weitergeleitet, auf welcher die Daten aus Testfall \RefFuncBrown{/T20-50/} eingetragen sind. Über die Schaltfläche 'Löschen' wird diese Kurseinheit entfernt und der Kursleiter wird auf die Seite 'Kursdetails' des Kurses 'Standardtanz' weitergeleitet. Auf dieser Seite ist nun im Bereich 'Kurseinheiten' die gelöschte Kurseinheit vom '24.08.2015' nicht mehr aufgeführt.(vgl. \RefFuncBlue{/F30-80/})			
		\end{itemize}
			
	\section{Testfälle für anonymer Benutzer}
			\stepHc
			\begin{itemize}
				\item \Test{Kursangebot ansehen} 
				 Der anonyme Benutzer klickt auf der Startseite auf den Button 'Unser Kursangebot' und gelangt auf die Seite 'Kursangebot durchsuchen'. Angezeigt wird das gesamte Kursangebot aufsteigend nach dem Beginn der nächsten Kurseinheit der Kurse, welche in den Testfällen \RefFuncBrown{/T10-80/}, \RefFuncBrown{/T10-90/} und in den Testfälle von \RefFuncBrown{/T20-10/} bis \RefFuncBrown{/T20-70/} angelegt wurden. (vgl. \RefFuncBlue{/F10-30/}) 
				 
				 \item \Test{Kurs suchen und Detailansicht ansehen} 
				 Der anonyme Benutzer befindet sich auf der Seite 'Kursangebot durchsuchen'. Angezeigt werden die Daten wie im Testfall \RefFuncBrown{/T30-10/}. 
					 \begin{itemize}
					 \item Der Benutzer wählt auf dieser Seite im Bereich 'Kursangebot durchsuchen' aus der Drop-Down Liste 'Kursname' aus und trägt im der darunterliegenden Eingabefeld den Namen 'Fußball' ein. Nach dem Drücken des Button 'Suchen' wird der Benutzer mit einer Fehlermeldung darauf hingewiesen, dass kein Kurs mit dem eingegebenen Namen existiert. Das entsprechende Eingabefeld wird geleert und der Benutzer bleibt auf Seite 'Kursangebot durchsuchen.
					 
					 \item  Der Benutzer wählt auf dieser Seite im Bereich 'Kursangebot durchsuchen' aus der Drop-Down Liste 'Kursname' aus und trägt im der darunterliegenden Eingabefeld den Namen 'Standardtanz' ein. Nach dem Klicken auf die Schaltfläche 'Suchen' wird in der Anzeige des Kursangebotes nur noch der Kurs 'Standardtanz' angezeigt.
					\end{itemize}
				(vgl. \RefFuncBlue{/F10-50/})	
				Anschließend klickt der Nutzer auf den angezeigten Kurs 'Standardtanz' und wird auf dessen Seite 'Kursdetails' weitergeleitet. Hier werden dem Benutzer die Identifikationsnummer, der Titel, die Beschreibung, der Kursleiter und die Kurseinheiten angezeigt, die in den Testfällen \RefFuncBrown{/T10-90/} und \RefFuncBrown{/T20-40/} bis \RefFuncBrown{/T20-70/} angelegt wurden. (vgl. \RefFuncBlue{/F10-60/}) 
				
					
				\item \Test{Registrierung} 
				 Der anonyme Benutzer befindet sich auf der Startseite. Über die Schaltfläche 'Registrieren' gelangt er auf die Seite 'Registrierung', welche das leere Registrierungsformular enthält. 		
					\begin{itemize}
						\item Der Benutzer registriert sich mit folgenden Daten: 
							\begin{itemize}
								\item Benutzername: 'Ricky1'
								\item Passwort: 'bSdF'
								\item Passwort: 'bSdF'
								\item Anrede: 'Frau'
								\item Name: 'Hölzl'
								\item Vorname: 'Katharina'
								\item Geburtsdatum: '29.05.1993'
								\item Adresse: 'Am Kastenfeld 39'
								\item Postleitzahl: '94081'
								\item Ort: 'Fürstenzell'
							\end{itemize}
						Der Benutzer drückt den Button 'Registrieren'. Daraufhin wird er durch verschiedene Fehlermeldungen darauf hingewiesen, dass der Benutzername 'Ricky1' bereits existiert, dass das Passwort zu kurz ist, dass keine E-Mailadresse angegeben wurde und dass die AGBs nicht bestätigt wurden. Die entsprechenden Felder Benutzername und Passwort werden geleert, alle anderen bereits eingegebenen Daten bleiben bestehen.
						
						\item Der Benutzer ändert seinen Benutzernamen auf 'Kathi5'. Außerdem gibt er die beiden Passwörter 'bSdFg7HjK' und 'bSdFg7HjK1' ein. Als E-Mailadresse gibt er 'katharinahoelzlweb.de' an. Nach Drücken des Button 'Registrieren' wird der Benutzer durch Fehlermeldungen darauf hingewiesen, dass seine Passwörter nicht übereinstimmen, die E-Mailadresse ein falsches Format hat und die AGBs nicht bestätigt wurden. Die Eingabefelder der Passwörter und der E-Mailadresse werden gelöscht, alle anderen bereits eingegebenen Daten bleiben bestehen.
						
						\item der Benutzer gibt ergänzt nun folgende Daten: 
							\begin{itemize}

								\item Passwort: 'bSdFg7HjK8*'
								\item Passwort: 'bSdFg7HjK8*'

								\item E-Mailadresse: 'katharina\_hoelzl@web.de'
							\end{itemize}
						Außerdem bestätigt er die AGBs durch Setzen des entsprechenden Hackens. Nach dem Drücken des Buttons 'Registrieren' wird das Formular abgeschickt und der Benutzer bekommt an seine angegebene E-Mailadresse einen Verifizierungslink geschickt. Durch die Eingabe dieses Links im Browser wird die E-Mailadresse bestätigt und der Benutzer wird anschließend auf die Startseite weitergeleitet.			
				\end{itemize}	
			(vgl. \RefFuncBlue{/F10-20/})
		\end{itemize}
		
	\section{Testfälle für den registrierten Benutzer}
		\stepHc
		\begin{itemize}
			\item \Test{Anmeldeversuch ohne vorherige Accountbestätigung}
			Der Benutzer befindet sich auf der Startseite, nachdem er sich registriert hat. (vgl. \RefFuncBrown{/T30-30/}) Nach Drücken des Button 'Anmelden' öffnet sich ein Drop-Down Menü, in welches der Benutzer seinen Benutzernamen 'Kathi5' und sein Passwort 'bSdFg7HjK' eingibt. Nach dem Klicken auf die Schaltfläche 'Anmelden' bekommt der Benutzer eine Fehlermeldung, welche ihn darauf hinweist, dass sein Account noch nicht bestätigt wurde. Der Benutzer wird auf die Startseite weitergeleitet.
			
			\item \Test{Account bestätigen durch den Systemadministrator}
			Der Systemadministrator befindet sich auf der Seite 'Adminverwaltung' und wird nach dem Drücken des Button 'Benutzer verwalten' auf die Seite 'Benutzer durchsuchen' weitergeleitet. Hier werden die bereits registrierten Benutzer 'Basti3', 'Ricky1' und 'Kathi5' mit den in den Testfällen \RefFuncBrown{/T10-50/}, \RefFuncBrown{/T10-60/} und \RefFuncBrown{/T30-30/} angelegten Daten angezeigt. Auf dieser Seite wählt der Admin im Drop-Down Menü 'Nicht aktiviert' aus und klickt anschließend auf den Button 'suchen'. Nun wird nur noch der Benutzer 'Kathi5' angezeigt. Der Administrator aktiviert den Account von 'Kathi5', indem er in das Kästchen rechts außen einen Hacken setzt und anschließend auf den Button 'Benutzer aktivieren' klickt. Der Benutzer 'Kathi5' ist nun aktiviert und kann sich ins System einloggen. (vgl. \RefFuncBrown{/T40-30/})  
			(vgl. \RefFuncBlue{/F40-110/})
			
			
			\item \Test{Einloggen ins System} 
			Der Benutzer 'Kathi5' klickt auf der Startseite	auf den 'Anmelden'-Button und es öffnet sich ein Drop-Down Menü, in welchem er Benutzername und Passwort eingeben kann. (vgl. \RefFuncBlue{/F20-10/})
			\begin{itemize}
				\item Der Benutzer gibt seinen Benutzernamen 'Kathi5' und das falsches Passwort 'wertz6uio' ein. Nach Drücken des 'Anmelden'-Buttons erhält er die Fehlermeldung 'Ihr Benutzername oder Ihr Passwort ist falsch'.

				\item Der Benutzer gibt den falschen Benutzernamen 'Kathi1' und das richtige Passwort 'bSdFg7HjK8*' ein. Nach Drücken des 'Anmelden'-Buttons erhält er die Fehlermeldung 'Ihr Benutzername oder Ihr Passwort ist falsch'.
				\item Der Benutzer gibt seinen Benutzernamen 'Kathi5' und kein Passwort ein. Nach Drücken des 'Anmelden'-Buttons erhält er die Fehlermeldung 'Ihr Benutzername oder Ihr Passwort ist falsch'.
				\item Der Benutzer gibt seinen Benutzernamen 'Kathi5' und das richtige Passwort 'bSdFg7HjK8*' ein. Nach Drücken des 'Anmelden'-Buttons wird der Benutzer auf die Seite 'Meine Kurse' weitergeleitet. Es sind hier noch keine Kurse eingetragen.

			\end{itemize} 
			
			\item \Test{Einloggen ins System} 
			Der Benutzer 'PatrickC' klickt auf der Startseite	auf den 'Anmelden'-Button und es öffnet sich ein Drop-Down Menü, in welchem er seinen Benutzernamen 'PatrickC' und sein Passwort 'lKjHg789' eingibt. Durch Drücken der Schaltfläche 'Anmelden' im Drop-Down Menü wird der Benutzer eingeloggt und auf die Seite 'Meine Kurse' weitergeleitet. Dort sind im Moment keine Kurse eingetragen. (vgl. \RefFuncBlue{/F20-10/})
			
			\item \Test{Einloggen ins System} 
			Der Benutzer 'Ricky1' klickt auf der Startseite	auf den 'Anmelden'-Button und es öffnet sich ein Drop-Down Menü, in welchem er seinen Benutzernamen 'Ricky1' und sein Passwort 'vBnMk89po' eingibt. Durch Drücken der Schaltfläche 'Anmelden' im Drop-Down Menü wird der Benutzer eingeloggt und auf die Seite 'Meine Kurse' weitergeleitet. Dort sind im Moment keine Kurse eingetragen. (vgl. \RefFuncBlue{/F20-10/})
				
			\item \Test{Konto aufladen} 
			Der Benutzer 'Kathi5' klickt in seinem Profil auf die Schaltfläche 'Konto aufladen' und wird auf die Seite 'Kontoaufladung' weitergeleitet. 
				\begin{itemize}
					\item Hier gibt Kathi5 folgende Daten ein:
						 \begin{itemize}
						 	\item Benutzerid: 'autogenerierte ID'
						 	\item Benutzername: 'Kathi5'
						 	\item Name: 'Hölzl'
						 	\item Vorname: 'Katharina'
						 	\item Kreditinstitut: ???
						 	\item Kreditkartennummer: ???
						 	\item Betrag: 'fünfzig'
						 \end{itemize}
					Anschließend klickt der Nutzer auf den Button 'Konto aufladen' und wird durch eine Fehlermeldung darauf hingewiesen, dass ein falsches Format im Eingabefeld 'Betrag' eingetragen wurde. Dieses Feld wird geleert, die anderen Eingaben bleiben erhalten.
						
					\item 'Kathi5' trägt in das Eingabefeld für den Betrag nun '50,00' ein. Nach dem Drücken des Button 'Aufladen' wird der Nutzer auf seine Profilseite weitergeleitet. Der Kontostand, welcher zuvor 0,00 Euro betrug, zeigt nun 50,00 an.  (vgl. \RefFuncBlue{/F20-180/})
				\end{itemize}
			\end{itemize}
			
			\begin{itemize}
				\item \Test{Kurs anmelden} 
					\begin{itemize}
						\item Der Benutzer 'Kathi5' befindet sich auf der Seite 'Kursangebot durchsuchen'. Es werden die Kurse wie im Testfall \RefFuncBrown{/T30-10/} angezeigt. Der Benutzer wählt im Bereich 'Kursangebot durchsuchen' aus der Drop-Down Liste 'Kostenlos' aus. Nach Bestätigen der Auswahl durch Drücken des Button 'Suchen', wird der Kurs 'Yoga' angezeigt. (vgl. \RefFuncBlue{/F20-160/})
						
						\item Der Benutzer wählt den Kurs 'Yoga' aus und wird auf die Seite 'Kursdetails' weitergeleitet. Angezeigt werden die Daten, die in den Testfällen \RefFuncBrown{/T10-80/} und \RefFuncBrown{/T20-10/} angelegt wurden. 'Kathi5' setzt einen Hacken in dem Kästchen 'Kursnews erhalten', um Änderungen im Kurs per E-Mail zu empfangen. Anschließend klickt sie auf den Button 'Anmelden' und wird auf die Seite 'Meine Kurse' weitergeleitet. Hier ist nun der Yoga-Kurs eingetragen und wird angezeigt. (vgl. \RefFuncBlue{/F20-120/})
					\end{itemize}
				
				\item \Test{Kurseinheiten anmelden} 
				Der Nutzer 'Kathi5'	klickt auf der Seite 'Meine Kurse' auf den Kurs 'Yoga' und erhält dessen Detailansicht. Es werden die Daten aus den Testfällen \RefFuncBrown{/T10-80/}  und \RefFuncBrown{/T20-10/} bis \RefFuncBrown{/T20-30/} angezeigt. Durch Klicken des Buttons 'Alle auswählen' und anschließend der Schaltfläche 'Speichern', trägt sich der Benutzer nun in alle verfügbaren Kurseinheiten ein. Anschließend wird er wieder auf die Seite 'Meine Kurse' zurückgeleitet. (vgl. \RefFuncBlue{/F20-190/})
				
				\item \Test{Kurseinheiten anmelden} 
				Der Benutzer 'PatrickC' befindet sich auf der Seite 'Kursangebot durchsuchen' und wählt den Kurs 'Standardtanz aus. Durch Anklicken dieses Kurses wird er auf die Detailseite weitergeleitet. Dort meldet er sich über die Schaltfläche 'Anmelden' für diesen Kurs an und wird auf die Seite 'Meine Kurse' weitergeleitet. Nun klickt 'PatrickC' auf den Kurs 'Standardtanz' und erhält dessen Detailansicht mit den Daten wie in den Testfällen \RefFuncBrown{/T20-40/} bis \RefFuncBrown{/T20-70/} festgelegt. Durch Setzen des Hackens bei dem Termin '04.08.2015' und dem anschließenden Drücken der Schaltfläche 'Speichern', trägt sich der Benutzer in die ausgewählte Kurseinheit ein. Anschließend wird er wieder auf die Seite 'Meine Kurse' zurückgeleitet. (vgl. \RefFuncBlue{/F20-190/})
				
				\item \Test{Kurs bezahlen}
				Der Benutzer 'Kathi5' meldet sich analog zum Testfall \RefFuncBrown{/T40-70/} für den Kurs 'Standardtanz' an. Durch Klicken auf den Kurs erhält der Nutzer die Detailansicht und wählt die folgenden Einheiten durch Setzen von Häckchen aus:
					\begin{itemize}
						\item 04.08.2015 
						\item 10.08.2015 
						\item 17.08.2015 
					\end{itemize}	
				Durch Drücken der Schaltfläche 'Speichern' trägt sich 'Kathi5' in die ausgewählten Kurseinheiten ein. Da eine Kurseinheit 15,00 Euro kostet, werden automatisch durch das Anmelden 45,00 Euro von ihrem Guthabenkonto abgebucht. Der Benutzer wird auf die Seite 'Meine Kurse' weitergeleitet. Der Kurs 'Standardtanz' wird nun angezeigt. (vgl. \RefFuncBlue{/F20-200/})
		 
				\item \Test{Kostenpflichtige Kurseinheit anmelden ohne genügend Guthaben}
				Der Benutzer 'Kathi5' befindet sich auf der Seite 'Meine Kurse' und gelangt durch Klicken auf den Kurs 'Standardtanz' zur dessen Detailansicht. Dort wählt er durch Setzen des Häckchens die Kurseinheiten '31.08.2015' und '07.09.2015' aus und bestätigt durch Drücken des Button 'Speichern'. Der Benutzer erhält eine Fehlermeldung, dass durch das Belegen dieser Kurseinheiten sein Guthaben über die erlaubte Grenze hinaus überschritten wird und er sich deshalb nicht anmelden kann. Die Häckchen der beiden ausgewählten Kurseinheiten werden entfernt, der Benutzer bleibt auf der Seite 'Detailansicht'.
				
				\item \Test{Kurseinheit abmelden und Rückbuchung des Kurspreises}
				Der Benutzer 'Kathi5' entfernt in der Detailansicht des Kurses 'Standardtanz' bei der Kurseinheit '17.08.2015' den Hacken im Kästchen 'Angemeldet und drückt den Button 'Speichern'. Dadurch wurde 'Kathi5' aus dieser Kurseinheit ausgetragen und der Betrag von 15,00 Euro wird automatisch auf das Guthabenkonto rückgebucht.  (vgl. \RefFuncBlue{/F20-200/}) und (vgl. \RefFuncBlue{/F20-210/})
				
				\item \Test{aktuelles Guthaben ansehen}
				Der Nutzer 'Kathi5' befindet sich auf der Profilseite. Hier wird der aktuelle Kontostand von 20,00 Euro angezeigt. (vgl. \RefFuncBlue{/F20-80/})
				
				\item \Test{Hilfeseite anzeigen}
				Der Benutzer 'Kathi5' befindet sich momentan auf der Profilseite und lässt sich durch Drücken des Buttons 'Hilfe' die Hilfeseite anzeigen. Dort sucht er die Information zum Thema 'Profilbild hochladen'. (vgl. \RefFuncBlue{/F10-80/})
					
				\item \Test{Hochladen eines Profilbildes}
				Der Nutzer 'Kathi5' befindet sich auf der Profilseite. Es werden die im Testfall \RefFuncBrown{/T30-30/} eingetragenen Daten angezeigt. Als Profilbild erscheint ein Dummy-Bild. 
					\begin{itemize}
						\item Zunächst wird der Button 'Bearbeiten' angeklickt. Dieser ändert sich in den Button 'Speichern'. Nun können die Angaben bearbeitet werden. Nach Drücken der Schaltfläche 'Durchsuchen' wird das Video 'Test.mp3' ausgewählt. Der Benutzer klickt auf den Button 'Speichern und es erscheint eine Fehlermeldung, da eine Musikdatei hochgeladen wurde. Das Dummy-Bild bleibt erhalten, ebenso wie der Button 'Speichern'.
						\item Nach erneutem Drücken der Schaltfläche 'Durchsuchen' wird das Bild 'KathiProfil.jpg' ausgewählt und Klicken auf den Button 'Speichern' hochgeladen. Es erscheint nun anstelle des Dummy-Bildes.	Der Button 'Speichern' wird wieder in 'Bearbeiten' geändert und die Eingabefelder können nicht geändert werden.	
					\end{itemize}
				(vgl. \RefFuncBlue{/F20-60/})
									
				\item \Test{E-Mailadresse ändern}
				'Kathi5' befindet sich auf der Profilseite. Nach Drücken der Schaltfläche 'Bearbeiten' können die Daten, bis auf die Benutzerrolle, geändert werden. und der Button 'Bearbeiten' ändert sich in 'Speichern'. Der Benutzer ändert die E-Mailadresse auf 'hoelzlka@fim.uni-passau.de'. Anschließend klickt 'Kathi5' auf den Button 'Speichern' und erhält an die angegebene Mailadresse einen Verifizierungslink. Nach Bestätigung durch diesen Link wird der Benutzer auf die Seite 'Mein Profil' weitergeleitet, in dem die E-Mailadresse auf die Aktuelle geändert ist. (vgl. \RefFuncBlue{/F20-60/})
				
				\item \Test{Anmelden in Kurs der bereits voll ist.}
				Der Benutzer 'Ricky1' befindet sich auf der Seite 'Kursangebot durchsuchen' und wählt den Kurs 'Standardtanz aus. Durch Anklicken dieses Kurses wird er auf die Detailseite weitergeleitet. Dort meldet er sich über die Schaltfläche 'Anmelden' für diesen Kurs an und wird auf die Seite 'Meine Kurse' weitergeleitet. Nun klickt 'Ricky1' auf den Kurs 'Standardtanz' und erhält dessen Detailansicht. Durch Setzen des Hackens bei dem Termin '04.08.2015' und dem anschließenden Drücken der Schaltfläche 'Speichern', will sich der Benutzer in die ausgewählte Kurseinheit eintragen. Da die maximale Teilnehmerzahl aber auf 2 beschränkt ist und bereits zwei Benutzer eingetragen sind, erhält 'Ricky1' eine Fehlermeldung und bleibt auf der Seite 'Detailansicht'.
				
				\item \Test{Kursteilnehmer ansehen}
				Der Benutzer 'Ricky1' befindet sich auf der Detailseite des Kurses 'Standardtanz' und klickt auf die Schaltfläche 'Zu den Teilnehmern'. Dadurch wird er auf die Seite 'Teilnehmer' weitergeleitet, auf der der Benutzername und das Profilbild der beiden Teilnehmer 'Kathi5' und 'PatrickC' eingetragen sind. (vgl. \RefFuncBlue{/F20-40/})
				
				\item \Test{Terminplaner ansehen}
				Der Benutzer 'Kathi5' wählt die Registerkarte 'Terminplaner' aus. Hier werden die in den Testfällen \RefFuncBrown{/T40-70/}, \RefFuncBrown{/T40-80/}, \RefFuncBrown{/T40-100/} und \RefFuncBrown{/T40-120/} eingetragenen Kurse angezeigt. (vgl. \RefFuncGreen{/F20-230W/})
				
				\item \Test{Aus dem System abmelden}
				Der Benutzer 'Kathi5' befindet sich momentan auf der Profilseite. Durch Drücken der Schaltfläche 'Abmelden' wird 'Kathi5' aus dem System abgemeldet und zur Startseite weitergeleitet. (vgl. \RefFuncBlue{/F20-90/})
					
				\item \Test{Passwort vergessen}
				Der Benutzer 'Kathi5' will sich anmelden, hat aber das sein Passwort vergessen. 'Kathi5' Klickt auf den Button 'Anmelden' und anschließend in dem sich öffnenden Drop-Down Menü auf 'Passwort vergessen'. In dem nun erscheinenden Eingabefeld gibt der Benutzer seine im System gespeicherte E-Mailadresse 'hoelzlka@fim.uni-passau.de' (vgl. \RefFuncBrown{/F40-160/}) ein. An diese Adresse wird nun ein automatisch generiertes Passwort geschickt, welches der Benutzer nun zusammen mit seinem Benutzernamen in die entsprechenden Felder im Drop-Down Menü auf der Startseite eingibt. Durch Drücken des Button 'Anmelden' wird der Benutzer eingeloggt. (vgl. \RefFuncBlue{/F20-100/})
					
			\end{itemize}
\section{Testfälle mit Datensatz}
			\stepHc
			\begin{itemize}
				\item \Test{Benachrichtigung durch Kursleiter}
				'??? muss noch abgeklärt werden ???' (vgl. \RefFuncBlue{/F30-110/})
				
				\item \Test{Guthabenkonto aufladen durch Systemadministrator}
				Der Systemadministrator trägt auf der Seite 'Adminverwaltung' im Bereich 'Konto aufladen' folgende Daten ein:
				\begin{itemize}
					\item BenutzerId: 'autogen'
					\item Name: Strohmeier
					\item Betrag: 40,00
				\end{itemize}
				Durch Drücken des Buttons 'Aufladen' bestätigt der Systemadministrator die Kontoaufladung und bleibt auf der Seite 'Adminverwaltung' (vgl. \RefFuncBlue{/F40-130/})
				
				\item \Test{Benutzer zu Kurseinheit hinzufügen durch Kursleiter}
				Der Kursleiter 'Basti3' befindet sich auf der Seite 'Kurseinheiten bearbeiten' der Kurseinheit '11.05.2015' des Kurses 'Standardtanz' und trägt im Bereich 'Teilnehmer hinzufügen' den Namen 'Strohmeier', den Vornamen 'Ricky' und die BenutzerId ein. Durch Drücken des Button 'Hinzufügen' wird der Benutzer 'Ricky1' zu dieser Kurseinheit eingetragen und der Kursleiter auf die Seite 'Meine Kurse' weitergeleitet. (vgl. \RefFuncBlue{/F30-60/})
						
				\item \Test{Benutzerrolle aufwerten durch Systemadministrator} 
				Der Systemadministrator klickt auf der Seite 'Adminverwaltung' auf die Schaltfläche 'Benutzer verwalten' und wird auf die Seite 'Benutzer durchsuchen'. Durch Anklicken des Benutzers 'Kathi5' wird der Admin auf die Seite 'Benutzer bearbeiten' weitergeleitet. Hier klickt er auf den Button 'Bearbeiten', welcher sich zu 'Speichern' ändert und wählt im Status des Benutzers 'Kathi5' die neue Rolle 'Kursleiter' aus. Durch Drücken des Buttons 'Speichern' bestätigt der Administrator diese Änderung und wird auf die Seite 'Adminverwaltung' weitergeleitet.(vgl. \RefFuncBlue{/F40-90/})
				
				\item \Test{Teilnehmer aus Kurs entfernen durch Kursleiter}
				Der Kursleiter 'Basti3' befindet sich auf der Seite 'Detailansicht des Kurses 'Yoga'. Hier drückt er den Button 'Zu den Teilnehmern' und wird auf die Teilnehmerseite weitergeleitet. Es ist bis jetz nur der Teilnehmer 'Kathi5' eingetragen. Auf dieser Seite wählt der Kursleiter den Teilnehmer 'Kathi5' aus und entfernt diesen aus dem Kurs durch Setzten des Häckchens und durch Drücken der Schaltfläche 'Löschen'. (vgl. \RefFuncBlue{/F30-90/})
				
				\item \Test{Kursleiter zu Kurs hinzufügen durch Systemadministrator}
				Auf der Oberfläche 'Adminverwaltung' drückt der Administrator auf den Button 'Kurse verwalten' und wird auf die seite 'Kursangebot durchsuchen' weitergeleitet. Dort wählt er den Kurs 'Yoga' aus und wird auf die Seite 'Kursdetails' weitergeleitet, auf der er Name und E-Mailadresse des neuen Kursleiters 'Kathi5' eingibt und die Änderung durch Drücken des Buttons 'Hinzufügen' abspeichert. Der Administrator wird auf die Seite 'Adminverwaltung' weitergeleitet. (vgl. \RefFuncBlue{/F40-30/})
				
				\item \Test{Statistiken anzeigen}
				Auf der Seite 'Adminverwaltung' klickt der Administrator auf die Schaltfläche 'Statistiken anzeigen' und wird auf diese weitergeleitet. Nach Auswahl des Parameters 'Tag' wird ein Säulendiagramm mit den Einnahmen der entsprechenden Woche angezeigt. (vgl. \RefFuncBlue{/F40-120/})
				
				\item \Test{Kurs löschen durch Systemadministrator}
				Der Systemadministrator befindet sich auf der Seite 'Detailansicht' des Kurses 'Yoga'. Durch Drücken des Buttons 'Kurs löschen' entfernt der Admin den Kurs 'Yoga' und wird auf die Seite 'Adminverwaltung' weitergeleitet. (vgl. \RefFuncBlue{/F40-40/})
				
				\item \Test{Benutzer löschen durch Systemadministrator}
				Der Systemadministrator drückt auf der Seite 'Adminverwaltung' den Button 'Benutzer verwalten' und wird auf die Seite 'Benutzer durchsuchen' weitergeleitet. Es werden die Benutzer 'Kathi5', 'Basti3', Ricky1' und 'PatrickC' angezeigt. Der Admin wählt in der Drop-Down Liste den Parameter 'Benutzername' aus und trägt 'Ricky1' in das darunterliegende Eingabefeld ein. Nun wird nur noch der Benutzer 'Ricky1' angezeigt. Diesen wählt der Administrator aus und wird auf dessen Profilseite weitergeleitet. Durch Drücken des Buttons 'Benutzer löschen' wird 'Ricky1' aus dem System entfernt. Der Administrator wird auf die Seite 'Adminverwaltung' zurückgeleitet. (vgl. \RefFuncBlue{/F40-80/})
				
				\item \Test{Kurs löschen durch Kursleiter}
				Der Kursleiter 'Basti3' befindet sich auf der Seite 'Detailansicht' des Kurses 'Standardtanz'. Durch Drücken des Buttons 'Kurs löschen' entfernt der Kursleiter den Kurs 'Standardtanz' und wird auf die Seite 'Meine Kurse' weitergeleitet. Diese enthält nun keine Einträge mehr. (vgl. \RefFuncBlue{/F30-30/})
				
			\end{itemize}





\chapter{Entwicklungsumgebung}
\begin{tiny}
		MB
	\end{tiny}\\

    \section{Software}
        \subsection{Betriebssysteme}
            \begin{itemize}
            	\item Windows 8.1
            	\item Windows 7
            	\item Linux Mint 17.1
            \end{itemize}	
            
        \subsection{Dokumentation}
            \begin{itemize}
            	\item Latex (u.a. Texmaker 4.1, TeXstudio 2.6.6)
            \end{itemize}
            
        \subsection{Datenbank}
           \begin{itemize}
           	\item PostgreSQL 9.1
           \end{itemize}
           
        \subsection{Webserver}
            \begin{itemize}
            	\item Tomcat 8
            \end{itemize}
            
        \subsection{Entwicklung}
           \begin{itemize}
           	\item Yatta UML Lab
           	\item Eclipse Luna 4.4.2
           	\item Eclipse Plugin: \\
	           	\begin{itemize}
	           		\item jUnit
	           		\item EGit
	           		\item Mylyn
	           		\item Checkstyle  
	           	\end{itemize} 
           \end{itemize}
           
        \subsection{sonstige Software}
            \begin{itemize}
            	\item DB-Visualizer 9.0.6
            \end{itemize}
           
    \section{Orgware}
        \begin{itemize}
        	\item GitHub
        	\item Dropbox
        \end{itemize}
        
    \section{Hardware privat}
        \begin{itemize}
        	\item Intel Core i5, 1.80 GHz, 8GB RAM
        	\item Intel Core i5, 2.4 GHz, 4GB RAM
        	\item Intel Core i7, 3.4 GHZ, 8GB RAM
        	\item Intel Core i5, 1.7 GHz, 4GB RAM
        \end{itemize}
            
        
\printnoidxglossaries
\end{document}